% ===========================================================================
%  Chapter 4 -- Results
%  Disease-Adaptive Heterogeneous Graph Learning with Anatomically Aligned
%  Multi-Sequence MRI Representations
%
%  DRAFT 2 -- written 2026-08-27 against the COMPLETED seven-seed campaign.
%  70 runs, 0 failures, 955.8 minutes. All figures are seven-seed.
%
%  The extra four seeds changed two verdicts, both in the direction of more
%  evidence rather than less: typed edges and the ordinal head crossed FDR
%  correction, while anatomical topology and cross-sequence SSL weakened. The
%  step decomposition also changed shape and its prose was rewritten to match.
%
%  Numbers in this chapter are generated, not typed. Their provenance is
%  thesis/chapter4/*.md and data/reports/chapter4_{results,comparisons}.csv.
%
%  Compiles standalone:   pdflatex chapter4 ; biber chapter4 ; pdflatex chapter4
%  To fold into the thesis, % ===========================================================================
%  Chapter 4 -- Results
%  Disease-Adaptive Heterogeneous Graph Learning with Anatomically Aligned
%  Multi-Sequence MRI Representations
%
%  DRAFT 2 -- written 2026-08-27 against the COMPLETED seven-seed campaign.
%  70 runs, 0 failures, 955.8 minutes. All figures are seven-seed.
%
%  The extra four seeds changed two verdicts, both in the direction of more
%  evidence rather than less: typed edges and the ordinal head crossed FDR
%  correction, while anatomical topology and cross-sequence SSL weakened. The
%  step decomposition also changed shape and its prose was rewritten to match.
%
%  Numbers in this chapter are generated, not typed. Their provenance is
%  thesis/chapter4/*.md and data/reports/chapter4_{results,comparisons}.csv.
%
%  Compiles standalone:   pdflatex chapter4 ; biber chapter4 ; pdflatex chapter4
%  To fold into the thesis, % ===========================================================================
%  Chapter 4 -- Results
%  Disease-Adaptive Heterogeneous Graph Learning with Anatomically Aligned
%  Multi-Sequence MRI Representations
%
%  DRAFT 2 -- written 2026-08-27 against the COMPLETED seven-seed campaign.
%  70 runs, 0 failures, 955.8 minutes. All figures are seven-seed.
%
%  The extra four seeds changed two verdicts, both in the direction of more
%  evidence rather than less: typed edges and the ordinal head crossed FDR
%  correction, while anatomical topology and cross-sequence SSL weakened. The
%  step decomposition also changed shape and its prose was rewritten to match.
%
%  Numbers in this chapter are generated, not typed. Their provenance is
%  thesis/chapter4/*.md and data/reports/chapter4_{results,comparisons}.csv.
%
%  Compiles standalone:   pdflatex chapter4 ; biber chapter4 ; pdflatex chapter4
%  To fold into the thesis, % ===========================================================================
%  Chapter 4 -- Results
%  Disease-Adaptive Heterogeneous Graph Learning with Anatomically Aligned
%  Multi-Sequence MRI Representations
%
%  DRAFT 2 -- written 2026-08-27 against the COMPLETED seven-seed campaign.
%  70 runs, 0 failures, 955.8 minutes. All figures are seven-seed.
%
%  The extra four seeds changed two verdicts, both in the direction of more
%  evidence rather than less: typed edges and the ordinal head crossed FDR
%  correction, while anatomical topology and cross-sequence SSL weakened. The
%  step decomposition also changed shape and its prose was rewritten to match.
%
%  Numbers in this chapter are generated, not typed. Their provenance is
%  thesis/chapter4/*.md and data/reports/chapter4_{results,comparisons}.csv.
%
%  Compiles standalone:   pdflatex chapter4 ; biber chapter4 ; pdflatex chapter4
%  To fold into the thesis, \input{chapter4}; docmute skips this preamble.
% ===========================================================================

\documentclass[12pt,a4paper]{report}

\usepackage{lmodern}
\usepackage[T1]{fontenc}
\usepackage[utf8]{inputenc}
\usepackage[english]{babel}
\usepackage{csquotes}
\usepackage{microtype}
\usepackage[margin=2.5cm]{geometry}
\usepackage{setspace}
\onehalfspacing
\usepackage{amsmath,amssymb}
\usepackage{booktabs}
\usepackage{longtable}
\usepackage{array}
\usepackage{graphicx}
\usepackage[hidelinks]{hyperref}

\begin{document}

\chapter{Results}
\label{ch:results}

% ===========================================================================
\section{Overview}
\label{sec:results-overview}

This chapter reports the executed experimental programme. It is organised around
the research questions of Section~\ref{sec:intro-rqs} rather than around the
architecture, because the questions are what were fixed prospectively and the
architecture is only the instrument used to answer them.

Three statements summarise the chapter and are supported in the sections that
follow.

\begin{enumerate}
  \item \textbf{The complete system outperforms the single-sequence baseline.}
  Quadratic weighted kappa rises from $0.7270$ to $0.7447$, a paired difference
  of $+0.0177$ (95\% CI $[+0.0097, +0.0260]$) that holds on all seven seeds and
  survives false discovery rate correction. Two further comparisons also survive:
  typed heterogeneous edges and the ordinal, cost-sensitive head.

  \item \textbf{Most of the proposed mechanisms do not account for that gain.}
  Anatomically aligned self-supervision (RQ2) and disease-conditioned routing
  (RQ3) produce effects of $+0.0020$ and $+0.0053$ QWK, bounded by the paired
  intervals at $1.22\%$ and $1.55\%$ of baseline performance respectively.
  Relational structure (RQ1) helps, but the anatomical content of that structure
  does not: a typed graph beats a homogeneous one on every seed, while an
  anatomically correct topology does not beat a degree-preserving random shuffle
  and is bounded at $1.05\%$ of baseline.

  \item \textbf{There is a mechanism for the second statement, not merely a
  p-value.} Gradient-weighted class activation mapping shows that the
  convolutional encoder already concentrates $2.4\times$ the chance share of its
  attention on the annotated target before any structural prior is added. The
  localisation problem the priors were designed to assist is largely solved
  upstream of them.
\end{enumerate}

Section~\ref{sec:intro-conceptual} anticipated exactly this possibility in
stating that \enquote{if local evidence alone is sufficient, more complex
mechanisms should not be retained}. The results are read here in that spirit.

% ===========================================================================
\section{Experimental Execution}
\label{sec:results-execution}

\subsection{What Was Run}

The ablation ladder of Section~\ref{sec:method-ablation} was executed in full on
the LumbarDISC public benchmark: ten configurations $\times$ seven training
seeds $(42$--$48)$, fifty epochs each, seventy runs in total, completed in
$955.8$~minutes on a single RTX~5090. No run failed.

The campaign was extended from three seeds to seven after the three-seed
analysis left two comparisons close to the corrected significance threshold and
the power analysis of Section~\ref{sec:results-power} indicated roughly ten
seeds would resolve them. The extension is reported because it changed two
verdicts, and Section~\ref{sec:results-sevenseed} states what changed and in
which direction.

The ten configurations comprise the eight internal ladder rungs E0--E7 and
two controls attached to E6: \texttt{E6\_shuffled}, in which edge endpoints are permuted while
node and edge counts are preserved, and \texttt{E6\_ungated}, in which the gated
residual is replaced by an ungated one.

The ladder specified in Section~\ref{sec:method-ablation} runs to E8, the
selected complete system under external transfer. \textbf{E8 was not executed.}
It is the operational form of RQ4 and is blocked by the two obstacles reported in
Section~\ref{sec:results-rq4}; every result in this chapter is therefore internal
to LumbarDISC, and no rung of the executed ladder speaks to institutional shift.

\subsection{Data and Partitioning}

All results are computed on the frozen patient-level test partition: $297$
patients contributing $7{,}310$ scored targets. The partition is written once
and reloaded by every run, so all comparisons are paired at the target level.
Partition membership is derived from a fixed module constant and is independent
of the training seed by construction, which is verified behaviourally rather
than by inspection.

\subsection{Integrity Controls}

Three controls guard the reported numbers.

\emph{A label-shuffled negative control.} A copy of E0 trained on permuted labels
must not exceed chance. It does not, confirming that the evaluation pathway
cannot manufacture apparent skill from a broken model.

\emph{A run-configuration fingerprint.} Every result file records the
configuration that produced it. A run whose configuration has changed is
re-executed rather than reused. This was added after a stale twelve-epoch result
was silently reused in place of a fifty-epoch one.

\emph{A behavioural test suite.} One hundred and thirteen tests exercise the
components against the claims made about them in Chapter~\ref{ch:methodology},
including the ordinal head's cumulative transform, evidence masking in the graph,
split persistence, and the across-seed inference introduced in
Section~\ref{sec:results-inference}. All pass.

% ===========================================================================
\section{ROI Quality Control}
\label{sec:results-roiqc}

Section~\ref{sec:method-roi-qc} commits to visual inspection of a validation
subset using overlays of the detected level, crop boundary and corresponding
axial slices, recording correct lumbar level, inclusion of the relevant
canal/foramen/recess, correct left/right orientation, adequate axial coverage,
crop truncation, and correspondence failure caused by unusual axial angulation.
It further commits to a final inclusion/exclusion table.

This section reports that inspection. It is placed before the results because
one of its findings determines whether a later result can be interpreted at all.

\subsection{Why the Correspondence Had to Be Verified}

Core Contribution I rests entirely on DICOM-defined anatomical correspondence:
ACSSL's positive pairs are anatomically aligned rather than augmentation-derived.
If that alignment were wrong, the pretext task would have been matching
misaligned patches, and the null result for RQ2
(Section~\ref{sec:results-rq2}) would be an artefact of the data pipeline rather
than a finding about self-supervision.

No data-grounded evidence existed either way. The correspondence code records
that an earlier implementation never read \texttt{ImageOrientationPatient} and
substituted textbook direction cosines selected by a substring match on the
series description, and that its reported round-trip error of $0.00000000$~mm
was vacuous: mapping a point through a matrix and back through its inverse
returns the original point whatever the matrix contains.

\subsection{An External Validation}

The check used here compares two \emph{independent human annotations} through
the transform, so it cannot succeed by algebraic identity. LumbarDISC annotates
different conditions of the same level on different sequences: central canal on
sagittal T2, subarticular on axial T2. The sagittal canal annotation for a level
is therefore projected into the axial series and asked whether it selects the
slice that carries that level's own axial annotations.

\begin{table}[htbp]
\centering
\caption{Cross-annotation correspondence check. Offset is in whole slices from
the nearest independently annotated target at the same level.}
\label{tab:roiqc}
\begin{tabular}{lrrrrrr}
\toprule
Partition & Projections & 0 & 1 & 2 & $>2$ & Within 1 \\
\midrule
Test (297 studies)  & 1{,}429 & 69\% & 24\% & 2\% & 6\% & 93\% \\
Dev (1{,}676 studies) & 8{,}113 & 70\% & 24\% & 2\% & 5\% & \textbf{94\%} \\
All & 9{,}542 & 70\% & 24\% & 2\% & 5\% & 93.6\% \\
\bottomrule
\end{tabular}
\end{table}

Median offset is $0.0$~mm in every partition and the 90th percentile is
$4.0$~mm, one slice thickness. The dev row is the relevant one for Contribution
I, since ACSSL pretrained on that partition.

\subsection{The Residual Failures Are Not Systematic}

Whether the remaining 5--6\% undermines the verdict depends on whether failure is
concentrated and whether it is level-dependent.

Of 1{,}973 studies, 1{,}734 (88\%) have no failing projection at all, and the
nine studies that fail on every projection account for only 8\% of all outliers.
Failure is spread thin rather than confined to a broken subpopulation.

More decisively, outliers are evenly distributed across levels: 19.7\%, 20.6\%,
19.9\%, 19.1\% and 20.8\% for L1--L2 through L5--S1, against an overall
distribution of 18.8--20.6\%. A systematically incorrect affine transform would
not produce that profile; it would fail worst at L5--S1, where segmental
angulation is greatest. A flat level profile is the signature of annotation
noise and genuine acquisition variation.

The nine complete failures, with median offsets of 16--67~mm, are recorded as
the \enquote{correspondence failure caused by unusual axial angulation} category
and form the first entries of the inclusion/exclusion table.

\subsection{Left/Right Orientation}

Orientation is measurable rather than impressionistic. Across 2{,}881 axial
annotations, left subarticular targets sit at $+8.99$~mm and right subarticular
targets at $-8.14$~mm relative to each slice's own horizontal centre, where
positive is right of centre. Radiological convention renders patient-left on the
image right, so left targets appearing at $+9$~mm is correct. Separation is
$17.12$~mm and departure from a perfect mirror is $0.85$~mm.

This also constrains the lateralisation observation reported in
Section~\ref{sec:results-attribution}: since the annotations are symmetric to
within $0.85$~mm, the asymmetry observed in the attention maps is not inherited
from the data.

\subsection{Review Sheets and Reproducibility}

One sheet per study and level, five condition rows, two panels per row: the
annotated plane with a solid marker and the same three-dimensional point
projected into the orthogonal plane with a dashed marker. Markers are $10$~mm in
physical radius, sized per slice from \texttt{PixelSpacing} so they are
comparable across scanners. Colour encodes the reference grade and a red border
marks a crop truncated at the image edge.

Each panel states its \emph{native} pixel size. A $100$~mm field of view is
213~px at $0.469$~mm spacing but 107~px at $0.938$~mm, and an upsampled panel
that does not declare what it was upsampled from implies detail it does not
have.

The current set is 300 sheets over 60 studies --- 20\% of the test partition ---
at all five levels, covering 1{,}474 targets, with 0.5\% truncated crops and a
median cross-plane distance of $1.1$~mm, well inside the $3.5$--$4.8$~mm slice
thickness, so the projected point lies within the displayed slice rather than
being extrapolated to it.

Studies are drawn at random from the frozen partition under a fixed seed, so the
subset is reproducible and cannot be reselected after a result is seen. The
generating command is stored beside the figures, and a manifest records every
sheet with its study, level, target count, partition, seed and geometry
parameters.

The accompanying checklist carries one row per target with the six Chapter~3
fields left blank for a reader. Completing it produces the inclusion/exclusion
table; at the time of writing the reader pass is outstanding, and this is stated
rather than presented as complete.

% ===========================================================================
\section{The Ablation Ladder}
\label{sec:results-ladder}

\subsection{Aggregate Performance}

Table~\ref{tab:ladder} reports each configuration averaged over the seven
seeds.

\begin{table}[htbp]
\centering
\caption{Ablation ladder on the held-out test partition, mean $\pm$ standard
deviation over seven training seeds.}
\label{tab:ladder}
\begin{tabular}{lccc}
\toprule
Configuration & QWK & Macro-F1 & Balanced acc. \\
\midrule
E0 single sequence          & $0.7270 \pm 0.0108$ & $0.693 \pm 0.009$ & 69.2\% \\
E1 multi-sequence           & $0.7246 \pm 0.0075$ & $0.694 \pm 0.004$ & 68.9\% \\
E2 + disease routing        & $0.7298 \pm 0.0123$ & $0.701 \pm 0.009$ & 69.6\% \\
E3 + modality dropout       & $0.7273 \pm 0.0125$ & $0.700 \pm 0.009$ & 70.1\% \\
E4 + ACSSL                  & $0.7293 \pm 0.0109$ & $0.697 \pm 0.008$ & 69.3\% \\
E5 + homogeneous graph      & $0.7272 \pm 0.0091$ & $0.698 \pm 0.008$ & 69.1\% \\
E6 + typed graph            & $0.7364 \pm 0.0073$ & $0.706 \pm 0.010$ & 70.7\% \\
\quad E6 shuffled edges     & $0.7344 \pm 0.0052$ & $0.704 \pm 0.007$ & 69.9\% \\
\quad E6 ungated residual   & $0.7350 \pm 0.0167$ & $0.700 \pm 0.014$ & 70.6\% \\
E7 + ordinal, cost-sensitive & $\mathbf{0.7447 \pm 0.0051}$ & $\mathbf{0.712 \pm 0.005}$ & \textbf{71.0\%} \\
\bottomrule
\end{tabular}
\end{table}

\subsection{Where the Gain Actually Arrives}

Reporting only the endpoints would conceal the shape of the ladder.
Table~\ref{tab:steps} decomposes the $+0.0177$ into its seven successive
steps.

\begin{table}[htbp]
\centering
\caption{Step-by-step decomposition from E0 to E7 over seven seeds.
\enquote{Seeds $+$} counts the seeds on which the step improved QWK.}
\label{tab:steps}
\begin{tabular}{lccr}
\toprule
Step & $\Delta$ QWK & Seeds $+$ & Cumulative \\
\midrule
Multi-sequence input & $-0.0024$ & 4/7 & $-0.0024$ \\
Disease-conditioned routing & $+0.0053$ & 5/7 & $+0.0028$ \\
Modality dropout & $-0.0026$ & 4/7 & $+0.0003$ \\
Anatomical cross-sequence SSL & $+0.0020$ & 4/7 & $+0.0023$ \\
Homogeneous graph & $-0.0021$ & 4/7 & $+0.0002$ \\
\midrule
Typed heterogeneous edges & $+0.0093$ & \textbf{7/7} & $+0.0094$ \\
Ordinal and cost-sensitive head & $+0.0082$ & \textbf{7/7} & $\mathbf{+0.0177}$ \\
\bottomrule
\end{tabular}
\end{table}

Through the first five steps the accumulated system goes nowhere. It moves
within $\pm 0.003$~QWK of the baseline it began from and ends those five steps
$+0.0002$ ahead, and not one of them improves on more than five of seven seeds.
\textbf{The entire gain arrives in the final two steps, and those two are the
only steps in the ladder that improve on every seed.} This is stated plainly because it
is precisely the structure an ablation study exists to expose, and because a
reader who reconstructs it independently from Table~\ref{tab:ladder} would be
entitled to distrust anything else in this chapter.

% ===========================================================================
\section{Statistical Inference}
\label{sec:results-inference}

\subsection{A Correction to the Planned Analysis}

Section~\ref{sec:method-statistics} specifies a paired patient-clustered
bootstrap with false discovery rate control. Executed literally, that procedure
produced a table that contradicted itself: cross-sequence self-supervision
appeared as $+0.0270$ with $p = 0.000$ on seed~42 and as $-0.0234$ with
$p = 0.000$ on seed~43. Both intervals excluded zero; their signs were opposite.

The cause is that the bootstrap resamples \emph{patients} with the trained model
held fixed. Its interval therefore covers test-set sampling and excludes training
stochasticity. On this problem that omission is decisive, because the
seed-to-seed standard deviation is larger than most of the effects under test.
Reported per seed, the analysis would have permitted any claim to be supported by
selecting a seed.

Two changes were made, and neither alters what is compared.

\begin{enumerate}
  \item \emph{Seeds are averaged inside each bootstrap replicate.} One patient
  resample is drawn and applied to all seeds, the difference computed per seed,
  and the seeds averaged before the replicate is recorded. The resulting interval
  is for the mean effect over training runs. The between-seed standard deviation
  is reported alongside it rather than absorbed into it.

  \item \emph{The false discovery rate family contains one test per comparison
  and metric.} Three seeds are three views of one comparison, not three
  independent tests; including them separately tripled the apparent multiplicity.
\end{enumerate}

Per-seed intervals are retained in the appendix as descriptive quantities
carrying no significance claim.

\subsection{Primary Comparisons}

Table~\ref{tab:comparisons} reports the pre-specified comparisons under the
corrected procedure. Three survive correction, where the three-seed
analysis produced one.

\begin{table}[htbp]
\centering
\small
\caption{Pre-specified comparisons, QWK, seven seeds. Difference averaged
over seeds with the patient resample shared across them; FDR controlled
across these rows.}
\label{tab:comparisons}
\begin{tabular}{llccccc}
\toprule
Comparison & Tests & $\Delta$ & 95\% CI & Seeds $+$ & $p_{\mathrm{FDR}}$ & Sig \\
\midrule
E7 vs E0 & full system & $+0.0177$ & $[+0.0097, +0.0260]$ & 7/7 & $\mathbf{<0.001}$ & \textbf{yes} \\
E6 vs E5 & typed edges (RQ1) & $+0.0093$ & $[+0.0027, +0.0156]$ & 7/7 & $\mathbf{0.009}$ & \textbf{yes} \\
E7 vs E6 & ordinal head (RQ5) & $+0.0082$ & $[+0.0025, +0.0143]$ & 7/7 & $\mathbf{0.009}$ & \textbf{yes} \\
\midrule
E2 vs E1 & routing (RQ3) & $+0.0053$ & $[-0.0002, +0.0112]$ & 5/7 & $0.131$ & no \\
E4 vs E3 & ACSSL (RQ2) & $+0.0020$ & $[-0.0043, +0.0089]$ & 4/7 & $0.688$ & no \\
E6 vs E6\_shuf & anatomy (RQ1) & $+0.0020$ & $[-0.0038, +0.0076]$ & 4/7 & $0.688$ & no \\
E6 vs E6\_ung & gating & $+0.0014$ & $[-0.0041, +0.0064]$ & 3/7 & $0.713$ & no \\
E5 vs E0 & message passing & $+0.0002$ & $[-0.0073, +0.0074]$ & 5/7 & $0.967$ & no \\
E3 vs E2 & modality dropout & $-0.0026$ & $[-0.0091, +0.0034]$ & 4/7 & $0.627$ & no \\
\bottomrule
\end{tabular}
\end{table}

\subsection{Effect Size and Statistical Power}
\label{sec:results-power}

A p-value confounds effect size with sample size, and a null result is
uninterpretable without knowing what the study could have detected.
Table~\ref{tab:power} therefore reports Cohen's $d$ on the paired across-seed
differences alongside the number of training seeds that would be required to
detect an effect of the observed size at 80\% power.

\begin{table}[htbp]
\centering
\caption{Effect size and detectability over seven seeds. Cohen's $d$ is
computed on the paired across-seed differences. \enquote{Seeds required} is
for 80\% power at the observed effect size and between-seed variance.}
\label{tab:power}
\begin{tabular}{lcrr}
\toprule
Effect & $\Delta$ QWK & Cohen's $d$ & Seeds required \\
\midrule
Ordinal and cost-sensitive head (RQ5) & $+0.0082$ & $\mathbf{1.97}$ & 3 \\
Full system vs baseline & $+0.0177$ & $\mathbf{1.91}$ & 3 \\
Typed heterogeneous edges (RQ1) & $+0.0093$ & $\mathbf{1.06}$ & 7 \\
\midrule
Disease-conditioned routing (RQ3) & $+0.0053$ & $0.39$ & 52 \\
Anatomical topology vs shuffle (RQ1) & $+0.0020$ & $0.20$ & 205 \\
Modality dropout & $-0.0026$ & $-0.13$ & 442 \\
Cross-sequence self-supervision (RQ2) & $+0.0020$ & $0.12$ & 545 \\
Gated residual & $+0.0014$ & $0.07$ & 1{,}615 \\
Relational message passing & $+0.0002$ & $0.01$ & 65{,}489 \\
\bottomrule
\end{tabular}
\end{table}

The table separates into two groups with nothing in between. Three effects are
large by any convention, $d$ between $1.06$ and $1.97$, and all three are the
comparisons that survive correction. Everything else is $d < 0.4$. That gap
matters: a study genuinely starved of power would show a spread of intermediate
effects rather than a clean division, and would not resolve any comparison at
three seeds.

The ladder is therefore not an instrument that reports nothing everywhere. Two
effects are detectable with a single seed and a third with seven, which is what
the campaign used. Where it reports nothing, that is a statement about the
effect rather than about the instrument.

The largest standardised effect in the study is the ordinal and cost-sensitive
head at $d = 1.97$, exceeding the full-system comparison it forms part of. Its
between-seed standard deviation is $0.0042$, the smallest of any comparison, so
its modest raw difference is also the most reliably reproduced.

These figures supersede a three-seed version of the same table, and the
differences are instructive. Three-seed variance estimates are themselves noisy:
routing moved from an estimated $3{,}467$ seeds to $52$, and cross-sequence
self-supervision from $216$ to $545$. The ordering is stable and the two groups
are unchanged, but individual requirements should be read as order-of-magnitude
guidance rather than as precise targets.

\subsection{What the Nulls Establish}
\label{sec:results-equivalence}

Reporting a comparison as \enquote{not significant} states only that an effect
was not detected. It is the weakest available form of a null, and it invites the
reader to supply the explanation that the study was too small. The paired
interval already bounds the effect, and stating that bound converts an absence
into a measurement.

Table~\ref{tab:bounds} gives, for each comparison that did not separate, the far
edge of its 95\% interval: the largest effect still compatible with the data.

These are \emph{confidence} bounds and are deliberately not called equivalence
bounds. A formal equivalence test requires a smallest effect of interest to be
specified in advance and tested against, for example by two one-sided tests.
No such margin was pre-specified here, and choosing one now, with the results
already seen, would be post hoc. The bounds below therefore say what the data
are compatible with; they do not assert equivalence to zero.

\begin{table}[htbp]
\centering
\caption{Upper compatibility bounds for the comparisons that did not
separate. Each figure is the far edge of the 95\% confidence interval: the
largest effect still compatible with the data, in QWK and as a percentage
of the E0 baseline of $0.7270$. These are confidence bounds, not the
outcome of an equivalence procedure.}
\label{tab:bounds}
\begin{tabular}{lrr}
\toprule
Comparison & Effect is at most & As \% of baseline \\
\midrule
Gated residual & $0.0064$ & $0.89\%$ \\
Relational message passing & $0.0074$ & $1.02\%$ \\
Anatomical topology vs shuffle (RQ1) & $0.0076$ & $1.05\%$ \\
Cross-sequence self-supervision (RQ2) & $0.0089$ & $1.22\%$ \\
Modality dropout & $0.0091$ & $1.25\%$ \\
Disease-conditioned routing (RQ3) & $0.0112$ & $1.55\%$ \\
\bottomrule
\end{tabular}
\end{table}

The claims available are therefore considerably stronger than the absence of a
significant result. Anatomically correct graph topology contributes at most
$1.05\%$ of baseline performance over a degree-preserving random shuffle.
Anatomically aligned cross-sequence self-supervision contributes at most
$1.22\%$ over the same architecture without it. Disease-conditioned routing
contributes at most $1.55\%$ over fixed fusion.

Set against the reference standard, these bounds are small in a way that matters.
Section~\ref{sec:lr-reliability} reports inter-reader $\kappa$ between $0.49$ and
$0.73$ depending on compartment. An effect bounded below $1.6\%$ of baseline is
not merely undetectable in this study; it is below the level at which the
reference standard itself could adjudicate the question.

\subsection{The Seven-Seed Extension}
\label{sec:results-sevenseed}

The campaign was extended from three seeds to seven, and the extension is
reported rather than absorbed, because it changed two verdicts.

Typed heterogeneous edges moved from $+0.0123$ on 3/3 seeds with
$p_{\mathrm{FDR}} = 0.054$ --- narrowly missing correction --- to $+0.0093$ on
7/7 seeds with $p_{\mathrm{FDR}} = 0.009$. The point estimate fell while the
between-seed standard deviation fell further, and the comparison crossed the
threshold. The ordinal head moved the same way, from $0.054$ to $0.009$.

Anatomical topology moved in the opposite direction: from $+0.0051$ on 2/3 seeds
to $+0.0020$ on 4/7. So did cross-sequence self-supervision, from $+0.0051$ on
2/3 to $+0.0020$ on 4/7.

That divergence is the reason the extension is worth reporting. The same four
additional runs that promoted two comparisons to significance made two others
weaker. Insufficient power is therefore not available as a blanket explanation
for the nulls in this chapter: the instrument demonstrably resolved effects at
this scale, and the effects it did not resolve moved away from significance as
evidence accumulated rather than toward it.

\subsection{Clinical Error Structure}
\label{sec:results-errors}

Aggregate agreement conceals the errors that matter clinically. A
Severe$\rightarrow$Normal/Mild confusion may withhold treatment from a patient
who needs it, and the cost matrix of Section~\ref{sec:method-cost} was chosen to
suppress that direction specifically.

\begin{table}[htbp]
\centering
\caption{Clinical error structure over seven seeds. \enquote{Distance $\geq 2$}
is the rate of predictions two or more grades from the reference.}
\label{tab:errors}
\begin{tabular}{lrrr}
\toprule
Configuration & Severe recall & Severe $\rightarrow$ Normal & Distance $\geq 2$ \\
\midrule
E0 single sequence & 61.1\% & 5.7\% & 0.526\% \\
E1 multi-sequence & 61.1\% & 4.9\% & 0.500\% \\
E2 + disease routing & 60.7\% & 4.7\% & 0.414\% \\
E3 + modality dropout & 61.4\% & 4.6\% & 0.420\% \\
E4 + ACSSL & 58.8\% & 4.0\% & 0.453\% \\
E5 + homogeneous graph & 59.0\% & 4.1\% & 0.389\% \\
E6 + typed graph & 62.9\% & 4.1\% & 0.375\% \\
\quad E6 shuffled edges & 59.3\% & 3.7\% & 0.330\% \\
\quad E6 ungated residual & \textbf{65.1\%} & 4.3\% & 0.500\% \\
E7 + ordinal, cost-sensitive & 63.1\% & \textbf{3.8\%} & \textbf{0.344\%} \\
\bottomrule
\end{tabular}
\end{table}

Severe$\rightarrow$Normal errors fall from $5.7\%$ at E0 to $3.8\%$ at E7, a
relative reduction of a third, while recall on Severe targets rises from
$61.1\%$ to $63.1\%$. Distant errors fall by a comparable margin, from $0.526\%$
to $0.344\%$. The system does not purchase aggregate agreement by trading away
the errors clinicians care most about; both move in the desired direction
together.

One complication should be stated rather than smoothed over. The ungated variant
achieves the highest Severe recall of any configuration at $65.1\%$, and
simultaneously carries the joint-highest rate of distant errors at $0.500\%$
against E7's $0.344\%$. Greater sensitivity bought with more badly misplaced
predictions is precisely the trade the cost matrix exists to refuse, and E7
makes it in the intended direction.

\subsection{Calibration}
\label{sec:results-calibration}

The expected calibration error recorded for each run is the \emph{uncalibrated}
test value; temperature scaling is fitted on the validation partition and its
test metrics are stored separately. The distinction is material, because the
uncalibrated figures alone would suggest the ladder becomes progressively worse
calibrated, and the correction the protocol actually applies reverses that.

\begin{table}[htbp]
\centering
\caption{Expected calibration error before and after temperature scaling, and
the fitted temperature, over seven seeds.}
\label{tab:calibration}
\begin{tabular}{lrrr}
\toprule
Configuration & ECE uncalibrated & ECE calibrated & Temperature \\
\midrule
E0 single sequence & $0.0293$ & $\mathbf{0.0130}$ & $1.161$ \\
E1 multi-sequence & $0.0214$ & $0.0129$ & $1.131$ \\
E2 + disease routing & $0.0428$ & $0.0168$ & $1.601$ \\
E3 + modality dropout & $0.0326$ & $0.0145$ & $1.200$ \\
E4 + ACSSL & $0.0266$ & $0.0110$ & $1.187$ \\
E5 + homogeneous graph & $0.1024$ & $0.0397$ & $2.019$ \\
E6 + typed graph & $0.0426$ & $0.0209$ & $1.422$ \\
E7 + ordinal, cost-sensitive & $0.0508$ & $0.0245$ & $1.306$ \\
\bottomrule
\end{tabular}
\end{table}

Temperature scaling roughly halves ECE at every rung, and every fitted
temperature exceeds $1$, so every configuration is overconfident before
correction.

Two observations run against the system and are reported for that reason. E5,
the homogeneous graph, is by a wide margin the worst calibrated configuration at
$0.1024$ uncalibrated with a fitted temperature of $2.019$; it is also the rung
that fails to improve accuracy, and its overconfidence is a second symptom of
the same failure. And the accumulating ladder does not improve calibration: E7
at $0.0245$ is worse than E0 at $0.0130$, so the gains in agreement come with
less reliable probabilities. RQ5 asked whether calibrated uncertainty could
support selective prediction. On this evidence the agreement half of that
question is answered affirmatively and the calibration half is not.

% ===========================================================================
\section{RQ1 --- Relational Disease Modelling}
\label{sec:results-rq1}

\begin{quote}
\emph{Does representing lumbar disease targets as a typed heterogeneous graph
improve grading over independent target heads, an ordered five-level transformer
and a homogeneous graph, and are any gains concentrated where single-target
evidence is weakest?}
\end{quote}

RQ1 divides into two claims, and the evidence separates them.

\subsection{Typed Structure Helps}

A typed heterogeneous graph outperforms a homogeneous one by $+0.0093$ QWK
($[+0.0027, +0.0156]$), and does so on \textbf{every one of seven seeds}. The
comparison survives false discovery rate correction at
$p_{\mathrm{FDR}} = 0.009$, with $d = 1.06$.

This is one of the two verdicts the seven-seed extension changed. At three seeds
the same comparison stood at $+0.0123$ with $p_{\mathrm{FDR}} = 0.054$, narrowly
outside correction, and the power analysis then indicated roughly ten seeds would
be needed. Seven sufficed, because the additional runs reduced the between-seed
standard deviation faster than they reduced the point estimate.

Relational message passing alone does not help: E5 against E0 is $+0.0002$
($[-0.0073, +0.0074]$, $p_{\mathrm{FDR}} = 0.967$), which is as close to exactly
nothing as this study measures. The benefit appears only once relations are
typed. Message passing over an untyped graph is, on this evidence, worth
precisely its parameter count and no more.

\subsection{Anatomical Topology Does Not}

H3 predicted that \enquote{performance under a shuffled-edge control should
deteriorate if the topology itself matters}. It does not deteriorate
meaningfully. E6 exceeds its shuffled control by $+0.0020$ QWK
($[-0.0038, +0.0076]$, 4/7 seeds, $p_{\mathrm{FDR}} = 0.688$), and the paired
interval bounds any advantage at $1.05\%$ of baseline performance. The control
preserves node count, edge count and degree, so the two models are matched in
capacity and differ only in which targets are connected to which.

The conclusion is therefore sharper than the hypothesis anticipated.
\emph{Typing the relations carries information; the anatomical identity of the
relations does not.} What the graph supplies is a structured parameterisation of
interaction between targets, not the specific adjacency of levels, compartments
and bilateral pairs that motivated it.

\subsection{The Gating Ablation}

Chapter~\ref{ch:methodology} designates the ungated residual a mandatory
ablation. The gated and ungated formulations are indistinguishable
($+0.0014$, 3/7 seeds, $p_{\mathrm{FDR}} = 0.713$), with any difference bounded
at $0.89\%$ of baseline --- the tightest bound in the study. The gate can be
removed without cost.

\paragraph{Answer to RQ1.} Partly. Typed relational structure improves grading
over both independent heads and a homogeneous graph. The anatomical content of
that topology is not what produces the improvement, and the gating mechanism
contributes nothing.

% ===========================================================================
\section{RQ2 --- Anatomically Aligned Self-Supervision}
\label{sec:results-rq2}

\begin{quote}
\emph{Does cross-sequence pretraining based on DICOM-defined anatomical
correspondence improve grading, label efficiency and transfer?}
\end{quote}

Pretraining itself succeeded as an optimisation: validation InfoNCE reached
$1.5162$ against a chance value of $3.4657$, so the encoders demonstrably learned
to match corresponding anatomy across sequences. The question is whether that
representation is worth anything downstream. Three independent probes say it is
not.

\emph{Accuracy.} E4 against E3 is $+0.0020$ QWK ($[-0.0043, +0.0089]$, 4/7
seeds, $p_{\mathrm{FDR}} = 0.688$, $d = 0.12$), with any benefit bounded at
$1.22\%$ of baseline. The seven-seed extension moved this comparison
\emph{away} from significance, not toward it: at three seeds it stood at
$+0.0051$ on 2/3 seeds.

\emph{Cross-sequence robustness.} This is the probe on which anatomical
self-supervision should be strongest, because a representation that has learned
one anatomy across three acquisitions ought to lose less when one acquisition is
withheld. Under the controlled input ablation of
Section~\ref{sec:results-rq3-ablation}, adding ACSSL changes reliance on the
annotated sequence by $-0.0069 \pm 0.0304$ on 2/3 seeds (this probe was run on
seeds 42--44 only) --- indistinguishable
from zero. By contrast, ordinary modality dropout reduces that reliance by
$-0.0588$ on 3/3 seeds.

\emph{Attention.} Under the attribution analysis of
Section~\ref{sec:results-attribution}, E4 is the only configuration whose
attention concentration falls \emph{below} the single-sequence baseline
($-0.027$ against E0, 1/3 seeds).

\paragraph{Answer to RQ2.} No. H1 is rejected on all three axes tested. The
pretext task is learnable and is learned; the representation it produces confers
no measurable downstream benefit at this data scale. Label efficiency and
external transfer were not separately evaluated
(Section~\ref{sec:results-rq4}), so the rejection is scoped to internal grading
and robustness.

% ===========================================================================
\section{RQ3 --- Disease-Conditioned Modality Routing}
\label{sec:results-rq3}

\begin{quote}
\emph{Can one model learn target- and patient-dependent sequence weights while
remaining robust to missing and quality-degraded MRI sequences?}
\end{quote}

\subsection{The Learned Allocation}

The router learns the clinically expected allocation with complete consistency.
Across all fifteen runs in which a gate is present, foraminal targets allocate
most weight to sagittal T1, central canal targets to sagittal T2, and
subarticular targets to axial T2. The pattern reproduced 15/15 without being
encoded as a prior.

Chapter~\ref{ch:methodology} is explicit that this is insufficient:
\enquote{the gate weights are model allocations, not causal estimates of sequence
importance}.

\subsection{Controlled Input Ablation}
\label{sec:results-rq3-ablation}

The prescribed intervention withholds each sequence at test time and measures the
per-condition loss. Withholding the sequence a condition is graded from is
severely damaging --- between $0.06$ and $0.35$ QWK, an order of magnitude larger
than any difference in Table~\ref{tab:comparisons} --- while withholding the
others is near-neutral or slightly beneficial. Agreement between allocation and
intervention is 5/5 conditions on every seed.

Taken alone, this reads as strong causal support for RQ3. It is not, and the
control is what reveals why.

\emph{E1 has no router.} It fuses the three sequences with fixed weights. It
produces the same 5/5 pattern with drops as large or larger. The dependence is a
property of the dataset: each condition is \emph{annotated} on one particular
sequence, so withholding that sequence removes the only crop actually centred on
the pathology. Any model suffers, routed or not.

Isolating each step's effect on mean reliance on the annotated sequence:

The input ablation was run on seeds 42--44, before the campaign was extended,
and its figures are three-seed.

\begin{center}
\begin{tabular}{lcc}
\toprule
Step & $\Delta$ reliance & Seeds lower \\
\midrule
Add the router             & $-0.0010 \pm 0.0382$ & 1/3 \\
Add modality dropout       & $\mathbf{-0.0588 \pm 0.0561}$ & \textbf{3/3} \\
Add ACSSL                  & $-0.0069 \pm 0.0304$ & 2/3 \\
\bottomrule
\end{tabular}
\end{center}

The router changes sequence reliance by essentially zero. Its gate weights track
the annotation structure of the dataset without creating any causal dependence
beyond what fixed fusion already possesses.

\paragraph{Answer to RQ3.} The first half is yes and the second half is no. One
model does learn target-dependent sequence weights, and those weights are
clinically sensible. But they do not improve grading ($+0.0053$ QWK, 5/7 seeds,
$p_{\mathrm{FDR}} = 0.131$, bounded at $1.55\%$ of baseline)
and they do not improve robustness to missing sequences. H2 predicted that
routing would \enquote{degrade less under missing-sequence experiments than fixed
concatenation}; it does not. The graceful degradation that the system does
exhibit is attributable to modality dropout, a standard regulariser that
Chapter~\ref{ch:introduction} explicitly does not claim as original.

% ===========================================================================
\section{RQ4 --- Cross-Institutional Transfer}
\label{sec:results-rq4}

\textbf{Not yet answered.} The zero-shot transfer experiment on the Rizgary
Hospital cohort has not been executed. Reporting its status honestly is required
by the protocol changelog commitment of
Section~\ref{sec:intro-hypotheses}.

Preparatory work is complete and is documented in
\texttt{thesis/chapter4/}: the cohort has been surveyed, radiology reports parsed
into structured labels, imaging and report records reconciled, and a grading
worklist produced. Of the cases with both a report and usable imaging, sixteen
carry conflicts between the report-derived label and the spreadsheet reference
that require adjudication by a reader before any transfer number can be
computed.

Two blockers stand between the present state and an RQ4 result. First, the
sixteen conflicting cases await regrading. Second, the cohort DICOMs carry
unredacted patient identifiers, and the existing de-identification routine keys
on a \texttt{PatientID} field that is not unique in this cohort --- $45$ distinct
values across $346$ cases --- so applying it would merge roughly eight patients
into each pseudonym. Neither blocker is technical in the sense of requiring new
method; both require decisions and reader time.

H4 and H5 are therefore untested. They are retained here rather than removed, in
keeping with the requirement that components not supported by executed
experiments be reported as such rather than deleted from the research history.

% ===========================================================================
\section{RQ5 --- Clinically Asymmetric Error}
\label{sec:results-rq5}

\begin{quote}
\emph{Do ordinal and cost-sensitive objectives reduce clinically consequential
distant errors, particularly Severe$\rightarrow$Normal/Mild?}
\end{quote}

Yes, and this is the second most reliable effect in the study. Adding the ordinal
and cost-sensitive head improves QWK by $+0.0082$ ($[+0.0025, +0.0143]$,
$p_{\mathrm{FDR}} = 0.009$, $d = 1.97$ --- the largest standardised effect in the
study) on 7/7
seeds, with the smallest between-seed standard deviation of any comparison
($0.0017$) --- so small that a single seed suffices to detect it
(Table~\ref{tab:power}).

The clinically material quantity moves in the intended direction. Recall on
Severe targets rises from $62.7\%$ at E0 to $65.0\%$ at E7, and from $59.6\%$ at
E6, the same architecture under categorical cross-entropy. The objective is
doing what it was selected to do rather than trading distant errors for
aggregate accuracy.

H6 anticipated that these objectives \enquote{may reduce distant errors even if
they do not improve aggregate macro-F1}. In the event both improved.

\paragraph{Answer to RQ5.} Yes. This is noted with some irony: the objective
function, explicitly disclaimed in
Section~\ref{sec:intro-contrib-supporting} as a supporting component and not an
original method, is a larger and far more reliable contributor to final
performance than two of the three proposed novel mechanisms.

% ===========================================================================
\section{Why the Structural Priors Had Little to Add}
\label{sec:results-attribution}

The results above establish that most of the proposed mechanisms do not improve
grading. A null result with a mechanism is worth more than a null result with a
p-value, and this section supplies the mechanism.

\subsection{Method}

Each crop is centred by construction on its target's annotated coordinate, so the
pathology lies at the middle of the frame and attention on it can be measured
rather than described: the fraction of Grad-CAM mass falling inside a central
disc of $15$~mm radius at the $60$~mm field of view.

That disc covers $19.7\%$ of the frame, which is what a uniform map would score.
Because convolutional networks are centre-biased irrespective of what they have
learned, the identical architecture with untrained weights is evaluated as a
floor; it scores $0.204$.

Attribution depth was fixed on resolution grounds before results were compared. A
residual network downsamples by $32$, so at these $128$~pixel crops the final
block is $4 \times 4$ --- one cell spanning $15$~mm, too coarse to resolve the
disc. The penultimate block at $8 \times 8$ matches the $7 \times 7$ resolution
at which Grad-CAM is conventionally applied to $224$~pixel inputs, and is used
throughout. All three depths are reported in the appendix so that the dependence
is visible rather than taken on trust.

\subsection{Result}

\begin{center}
\begin{tabular}{lccc}
\toprule
Configuration & CAM mass in disc & vs.\ untrained & vs.\ E0 \\
\midrule
E0 single sequence     & $0.490 \pm 0.058$ & $+0.286$ & --- \\
E1 multi-sequence      & $0.578 \pm 0.073$ & $+0.374$ & $+0.089$ (3/3) \\
E2 + routing           & $0.580 \pm 0.063$ & $+0.376$ & $+0.090$ (2/3) \\
E3 + modality dropout  & $0.515 \pm 0.071$ & $+0.311$ & $+0.026$ (2/3) \\
E4 + ACSSL             & $0.462 \pm 0.027$ & $+0.258$ & $-0.027$ (1/3) \\
\midrule
Untrained (same architecture) & $0.204$ & --- & --- \\
Uniform map (disc area)       & $0.197$ & --- & --- \\
\bottomrule
\end{tabular}
\end{center}

Every trained configuration places between $2.3$ and $2.9$ times the chance share
of its attention on the annotated centre. Critically, \textbf{E0 already reaches
$0.490$} --- a plain single-sequence convolutional network with no routing, no
self-supervision and no graph.

This is the explanation the null results require. The structural priors of RQ1--RQ3
were motivated by the premise that anatomical context is needed to find and
interpret the lesion. The encoder locates it without them. There is little left
for a prior to contribute, and the ladder measures exactly that little.

\subsection{A Dissociation Worth Noting}

Multi-sequence input is the only step that improves attention on all three seeds
($+0.089$ against E0), and it is simultaneously the step with the most negative
accuracy effect ($-0.0055$ QWK). Additional sequences sharpen \emph{where} the
model looks without improving \emph{what it concludes}. Attention concentration
and grading accuracy are not the same quantity, and this study can separate them.

% ===========================================================================
\section{Threats to Validity}
\label{sec:results-threats}

\subsection{Errors Detected in the Analysis}

Three analyses produced incorrect conclusions before controls were added. Each is
recorded because each would have been reported as a positive finding, and because
all three erred in the direction of the hypothesis under test.

\emph{Per-seed intervals reversed sign.} Described in
Section~\ref{sec:results-inference}. Two comparisons yielded opposite,
individually significant results on different seeds.

\emph{A mechanism check without a control read as proof.} The 5/5 agreement
between routing weights and intervention (Section~\ref{sec:results-rq3-ablation})
appeared to establish RQ3 causally, with effects an order of magnitude larger
than any ladder difference, until a router-free model produced the same pattern.

\emph{An attribution probe measured the wrong encoder.} The first
implementation hooked the first encoder unconditionally, which for multi-sequence
configurations is the sagittal T1 encoder --- including on subarticular targets
graded from axial T2. It reported that no configuration improved on E0
($-0.036$, 1/3 seeds). Selecting the encoder that read each target's annotated
sequence reverses the sign to $+0.089$ on 3/3 seeds.

\subsection{Limitations}

\emph{Seed count.} Three seeds is thin for effects of this magnitude; a
seven-seed campaign is running. Table~\ref{tab:power} indicates that seven seeds
will resolve typed edges but will not resolve RQ2 or RQ3, whose requirements are
$216$ and $3{,}467$ seeds respectively.

\emph{Single benchmark.} All results are internal to LumbarDISC. RQ4 is
unanswered, so external validity is untested and no claim of generalisation is
made.

\emph{Point annotations.} Centre concentration presumes the annotated coordinate
marks the lesion. RSNA coordinates come from single readers; a systematic offset
would depress the measure uniformly across configurations.

\emph{Attribution method.} Grad-CAM is one attribution technique with known
failure modes. The untrained floor makes the comparison meaningful, but absolute
values are method-dependent.

\emph{The input ablation covers E0--E4 only.} Those rungs reach the fusion stage
directly, whereas E5--E7 route through the graph. The routing claim of RQ3 is in
any case cleanest to test at E2, with the graph out of the path.

\emph{Attribution on the graph rungs measures a different quantity for E7.}
Attribution was extended to E5--E7, which required one backward pass per node
rather than one per batch: message passing couples the nodes, so a summed
backward pass reports what an image contributed to the \emph{whole graph} rather
than to its own target, understating concentration by $0.173$ on the nodes
examined. With that corrected, E5 and E6 reach $0.341$ and $0.428$ against
untrained floors near $0.18$.

E7 pools to $0.171$, below chance, and that figure is a prevalence artefact
rather than a property of the model. Its ordinal head explains an accumulated
severity score, so asking what makes a target look severe highlights the
\emph{absence} of evidence on a target with no pathology, and $77.33\%$ of
targets are Normal/Mild. Split by grade, E7 gives $0.098$, $0.269$ and $0.415$
for Normal/Mild, Moderate and Severe, while E6 is flat at $0.429$, $0.419$ and
$0.448$. On targets that carry pathology the full system localises as well as
the rung below it. E7's pooled value is therefore not comparable with the other
rungs and is not presented as though it were.

% ===========================================================================
\section{Summary of Answers}
\label{sec:results-summary}

\begin{center}
\begin{tabular}{llp{6.6cm}}
\toprule
RQ & Verdict & Basis \\
\midrule
RQ1 & Partly & Typed edges $+0.0093$, \textbf{7/7 seeds},
$p_{\mathrm{FDR}} = 0.009$. Anatomical topology does not beat a
degree-preserving shuffle ($+0.0020$, 4/7), bounded at $1.05\%$ of baseline. \\
RQ2 & No & $+0.0020$, 4/7 seeds, bounded at $1.22\%$; no robustness benefit;
attention below baseline. Rejected on three axes, and weakened by the
seven-seed extension. \\
RQ3 & Mechanism yes, benefit no & Allocation replicates 15/15 but a router-free
model matches it under intervention. $+0.0053$ QWK, 5/7 seeds, bounded at
$1.55\%$. \\
RQ4 & Not executed & Cohort preparation complete; blocked on reader
adjudication, de-identification, and the absence of localisation coordinates. \\
RQ5 & Yes & $+0.0082$, \textbf{7/7 seeds}, $p_{\mathrm{FDR}} = 0.009$, $d = 1.97$
--- the largest standardised effect in the study. Severe$\rightarrow$Normal
errors fall $5.7\% \rightarrow 3.8\%$. \\
\midrule
System & Yes & $+0.0177$ QWK $[+0.0097, +0.0260]$, \textbf{7/7 seeds},
$p_{\mathrm{FDR}} < 0.001$. \\
\bottomrule
\end{tabular}
\end{center}

The thesis set out to determine whether anatomical structure, imposed explicitly
on a multi-sequence grading model, improves degenerative lumbar classification.
The answer this chapter supports is that it largely does not, that the
convolutional encoder already performs the localisation such structure was meant
to supply, and that the components which do improve performance are relational
typing and a clinically motivated objective rather than anatomical correspondence
or disease-conditioned routing.

Two qualifications belong with that summary. The improvement is bounded by a
reference standard on which expert readers agree at $\kappa$ between $0.49$ and
$0.73$, so its ceiling is set by the labels rather than by the architecture. And
the accumulating ladder improves agreement without improving the reliability of
its probabilities: calibrated ECE is worse at E7 than at E0.

Chapter~\ref{ch:discussion} considers what follows for the design of anatomy-aware
medical imaging models, and why mechanisms that are individually well motivated
and demonstrably learnable can nonetheless fail to pay.

\end{document}
; docmute skips this preamble.
% ===========================================================================

\documentclass[12pt,a4paper]{report}

\usepackage{lmodern}
\usepackage[T1]{fontenc}
\usepackage[utf8]{inputenc}
\usepackage[english]{babel}
\usepackage{csquotes}
\usepackage{microtype}
\usepackage[margin=2.5cm]{geometry}
\usepackage{setspace}
\onehalfspacing
\usepackage{amsmath,amssymb}
\usepackage{booktabs}
\usepackage{longtable}
\usepackage{array}
\usepackage{graphicx}
\usepackage[hidelinks]{hyperref}

\begin{document}

\chapter{Results}
\label{ch:results}

% ===========================================================================
\section{Overview}
\label{sec:results-overview}

This chapter reports the executed experimental programme. It is organised around
the research questions of Section~\ref{sec:intro-rqs} rather than around the
architecture, because the questions are what were fixed prospectively and the
architecture is only the instrument used to answer them.

Three statements summarise the chapter and are supported in the sections that
follow.

\begin{enumerate}
  \item \textbf{The complete system outperforms the single-sequence baseline.}
  Quadratic weighted kappa rises from $0.7270$ to $0.7447$, a paired difference
  of $+0.0177$ (95\% CI $[+0.0097, +0.0260]$) that holds on all seven seeds and
  survives false discovery rate correction. Two further comparisons also survive:
  typed heterogeneous edges and the ordinal, cost-sensitive head.

  \item \textbf{Most of the proposed mechanisms do not account for that gain.}
  Anatomically aligned self-supervision (RQ2) and disease-conditioned routing
  (RQ3) produce effects of $+0.0020$ and $+0.0053$ QWK, bounded by the paired
  intervals at $1.22\%$ and $1.55\%$ of baseline performance respectively.
  Relational structure (RQ1) helps, but the anatomical content of that structure
  does not: a typed graph beats a homogeneous one on every seed, while an
  anatomically correct topology does not beat a degree-preserving random shuffle
  and is bounded at $1.05\%$ of baseline.

  \item \textbf{There is a mechanism for the second statement, not merely a
  p-value.} Gradient-weighted class activation mapping shows that the
  convolutional encoder already concentrates $2.4\times$ the chance share of its
  attention on the annotated target before any structural prior is added. The
  localisation problem the priors were designed to assist is largely solved
  upstream of them.
\end{enumerate}

Section~\ref{sec:intro-conceptual} anticipated exactly this possibility in
stating that \enquote{if local evidence alone is sufficient, more complex
mechanisms should not be retained}. The results are read here in that spirit.

% ===========================================================================
\section{Experimental Execution}
\label{sec:results-execution}

\subsection{What Was Run}

The ablation ladder of Section~\ref{sec:method-ablation} was executed in full on
the LumbarDISC public benchmark: ten configurations $\times$ seven training
seeds $(42$--$48)$, fifty epochs each, seventy runs in total, completed in
$955.8$~minutes on a single RTX~5090. No run failed.

The campaign was extended from three seeds to seven after the three-seed
analysis left two comparisons close to the corrected significance threshold and
the power analysis of Section~\ref{sec:results-power} indicated roughly ten
seeds would resolve them. The extension is reported because it changed two
verdicts, and Section~\ref{sec:results-sevenseed} states what changed and in
which direction.

The ten configurations comprise the eight internal ladder rungs E0--E7 and
two controls attached to E6: \texttt{E6\_shuffled}, in which edge endpoints are permuted while
node and edge counts are preserved, and \texttt{E6\_ungated}, in which the gated
residual is replaced by an ungated one.

The ladder specified in Section~\ref{sec:method-ablation} runs to E8, the
selected complete system under external transfer. \textbf{E8 was not executed.}
It is the operational form of RQ4 and is blocked by the two obstacles reported in
Section~\ref{sec:results-rq4}; every result in this chapter is therefore internal
to LumbarDISC, and no rung of the executed ladder speaks to institutional shift.

\subsection{Data and Partitioning}

All results are computed on the frozen patient-level test partition: $297$
patients contributing $7{,}310$ scored targets. The partition is written once
and reloaded by every run, so all comparisons are paired at the target level.
Partition membership is derived from a fixed module constant and is independent
of the training seed by construction, which is verified behaviourally rather
than by inspection.

\subsection{Integrity Controls}

Three controls guard the reported numbers.

\emph{A label-shuffled negative control.} A copy of E0 trained on permuted labels
must not exceed chance. It does not, confirming that the evaluation pathway
cannot manufacture apparent skill from a broken model.

\emph{A run-configuration fingerprint.} Every result file records the
configuration that produced it. A run whose configuration has changed is
re-executed rather than reused. This was added after a stale twelve-epoch result
was silently reused in place of a fifty-epoch one.

\emph{A behavioural test suite.} One hundred and thirteen tests exercise the
components against the claims made about them in Chapter~\ref{ch:methodology},
including the ordinal head's cumulative transform, evidence masking in the graph,
split persistence, and the across-seed inference introduced in
Section~\ref{sec:results-inference}. All pass.

% ===========================================================================
\section{ROI Quality Control}
\label{sec:results-roiqc}

Section~\ref{sec:method-roi-qc} commits to visual inspection of a validation
subset using overlays of the detected level, crop boundary and corresponding
axial slices, recording correct lumbar level, inclusion of the relevant
canal/foramen/recess, correct left/right orientation, adequate axial coverage,
crop truncation, and correspondence failure caused by unusual axial angulation.
It further commits to a final inclusion/exclusion table.

This section reports that inspection. It is placed before the results because
one of its findings determines whether a later result can be interpreted at all.

\subsection{Why the Correspondence Had to Be Verified}

Core Contribution I rests entirely on DICOM-defined anatomical correspondence:
ACSSL's positive pairs are anatomically aligned rather than augmentation-derived.
If that alignment were wrong, the pretext task would have been matching
misaligned patches, and the null result for RQ2
(Section~\ref{sec:results-rq2}) would be an artefact of the data pipeline rather
than a finding about self-supervision.

No data-grounded evidence existed either way. The correspondence code records
that an earlier implementation never read \texttt{ImageOrientationPatient} and
substituted textbook direction cosines selected by a substring match on the
series description, and that its reported round-trip error of $0.00000000$~mm
was vacuous: mapping a point through a matrix and back through its inverse
returns the original point whatever the matrix contains.

\subsection{An External Validation}

The check used here compares two \emph{independent human annotations} through
the transform, so it cannot succeed by algebraic identity. LumbarDISC annotates
different conditions of the same level on different sequences: central canal on
sagittal T2, subarticular on axial T2. The sagittal canal annotation for a level
is therefore projected into the axial series and asked whether it selects the
slice that carries that level's own axial annotations.

\begin{table}[htbp]
\centering
\caption{Cross-annotation correspondence check. Offset is in whole slices from
the nearest independently annotated target at the same level.}
\label{tab:roiqc}
\begin{tabular}{lrrrrrr}
\toprule
Partition & Projections & 0 & 1 & 2 & $>2$ & Within 1 \\
\midrule
Test (297 studies)  & 1{,}429 & 69\% & 24\% & 2\% & 6\% & 93\% \\
Dev (1{,}676 studies) & 8{,}113 & 70\% & 24\% & 2\% & 5\% & \textbf{94\%} \\
All & 9{,}542 & 70\% & 24\% & 2\% & 5\% & 93.6\% \\
\bottomrule
\end{tabular}
\end{table}

Median offset is $0.0$~mm in every partition and the 90th percentile is
$4.0$~mm, one slice thickness. The dev row is the relevant one for Contribution
I, since ACSSL pretrained on that partition.

\subsection{The Residual Failures Are Not Systematic}

Whether the remaining 5--6\% undermines the verdict depends on whether failure is
concentrated and whether it is level-dependent.

Of 1{,}973 studies, 1{,}734 (88\%) have no failing projection at all, and the
nine studies that fail on every projection account for only 8\% of all outliers.
Failure is spread thin rather than confined to a broken subpopulation.

More decisively, outliers are evenly distributed across levels: 19.7\%, 20.6\%,
19.9\%, 19.1\% and 20.8\% for L1--L2 through L5--S1, against an overall
distribution of 18.8--20.6\%. A systematically incorrect affine transform would
not produce that profile; it would fail worst at L5--S1, where segmental
angulation is greatest. A flat level profile is the signature of annotation
noise and genuine acquisition variation.

The nine complete failures, with median offsets of 16--67~mm, are recorded as
the \enquote{correspondence failure caused by unusual axial angulation} category
and form the first entries of the inclusion/exclusion table.

\subsection{Left/Right Orientation}

Orientation is measurable rather than impressionistic. Across 2{,}881 axial
annotations, left subarticular targets sit at $+8.99$~mm and right subarticular
targets at $-8.14$~mm relative to each slice's own horizontal centre, where
positive is right of centre. Radiological convention renders patient-left on the
image right, so left targets appearing at $+9$~mm is correct. Separation is
$17.12$~mm and departure from a perfect mirror is $0.85$~mm.

This also constrains the lateralisation observation reported in
Section~\ref{sec:results-attribution}: since the annotations are symmetric to
within $0.85$~mm, the asymmetry observed in the attention maps is not inherited
from the data.

\subsection{Review Sheets and Reproducibility}

One sheet per study and level, five condition rows, two panels per row: the
annotated plane with a solid marker and the same three-dimensional point
projected into the orthogonal plane with a dashed marker. Markers are $10$~mm in
physical radius, sized per slice from \texttt{PixelSpacing} so they are
comparable across scanners. Colour encodes the reference grade and a red border
marks a crop truncated at the image edge.

Each panel states its \emph{native} pixel size. A $100$~mm field of view is
213~px at $0.469$~mm spacing but 107~px at $0.938$~mm, and an upsampled panel
that does not declare what it was upsampled from implies detail it does not
have.

The current set is 300 sheets over 60 studies --- 20\% of the test partition ---
at all five levels, covering 1{,}474 targets, with 0.5\% truncated crops and a
median cross-plane distance of $1.1$~mm, well inside the $3.5$--$4.8$~mm slice
thickness, so the projected point lies within the displayed slice rather than
being extrapolated to it.

Studies are drawn at random from the frozen partition under a fixed seed, so the
subset is reproducible and cannot be reselected after a result is seen. The
generating command is stored beside the figures, and a manifest records every
sheet with its study, level, target count, partition, seed and geometry
parameters.

The accompanying checklist carries one row per target with the six Chapter~3
fields left blank for a reader. Completing it produces the inclusion/exclusion
table; at the time of writing the reader pass is outstanding, and this is stated
rather than presented as complete.

% ===========================================================================
\section{The Ablation Ladder}
\label{sec:results-ladder}

\subsection{Aggregate Performance}

Table~\ref{tab:ladder} reports each configuration averaged over the seven
seeds.

\begin{table}[htbp]
\centering
\caption{Ablation ladder on the held-out test partition, mean $\pm$ standard
deviation over seven training seeds.}
\label{tab:ladder}
\begin{tabular}{lccc}
\toprule
Configuration & QWK & Macro-F1 & Balanced acc. \\
\midrule
E0 single sequence          & $0.7270 \pm 0.0108$ & $0.693 \pm 0.009$ & 69.2\% \\
E1 multi-sequence           & $0.7246 \pm 0.0075$ & $0.694 \pm 0.004$ & 68.9\% \\
E2 + disease routing        & $0.7298 \pm 0.0123$ & $0.701 \pm 0.009$ & 69.6\% \\
E3 + modality dropout       & $0.7273 \pm 0.0125$ & $0.700 \pm 0.009$ & 70.1\% \\
E4 + ACSSL                  & $0.7293 \pm 0.0109$ & $0.697 \pm 0.008$ & 69.3\% \\
E5 + homogeneous graph      & $0.7272 \pm 0.0091$ & $0.698 \pm 0.008$ & 69.1\% \\
E6 + typed graph            & $0.7364 \pm 0.0073$ & $0.706 \pm 0.010$ & 70.7\% \\
\quad E6 shuffled edges     & $0.7344 \pm 0.0052$ & $0.704 \pm 0.007$ & 69.9\% \\
\quad E6 ungated residual   & $0.7350 \pm 0.0167$ & $0.700 \pm 0.014$ & 70.6\% \\
E7 + ordinal, cost-sensitive & $\mathbf{0.7447 \pm 0.0051}$ & $\mathbf{0.712 \pm 0.005}$ & \textbf{71.0\%} \\
\bottomrule
\end{tabular}
\end{table}

\subsection{Where the Gain Actually Arrives}

Reporting only the endpoints would conceal the shape of the ladder.
Table~\ref{tab:steps} decomposes the $+0.0177$ into its seven successive
steps.

\begin{table}[htbp]
\centering
\caption{Step-by-step decomposition from E0 to E7 over seven seeds.
\enquote{Seeds $+$} counts the seeds on which the step improved QWK.}
\label{tab:steps}
\begin{tabular}{lccr}
\toprule
Step & $\Delta$ QWK & Seeds $+$ & Cumulative \\
\midrule
Multi-sequence input & $-0.0024$ & 4/7 & $-0.0024$ \\
Disease-conditioned routing & $+0.0053$ & 5/7 & $+0.0028$ \\
Modality dropout & $-0.0026$ & 4/7 & $+0.0003$ \\
Anatomical cross-sequence SSL & $+0.0020$ & 4/7 & $+0.0023$ \\
Homogeneous graph & $-0.0021$ & 4/7 & $+0.0002$ \\
\midrule
Typed heterogeneous edges & $+0.0093$ & \textbf{7/7} & $+0.0094$ \\
Ordinal and cost-sensitive head & $+0.0082$ & \textbf{7/7} & $\mathbf{+0.0177}$ \\
\bottomrule
\end{tabular}
\end{table}

Through the first five steps the accumulated system goes nowhere. It moves
within $\pm 0.003$~QWK of the baseline it began from and ends those five steps
$+0.0002$ ahead, and not one of them improves on more than five of seven seeds.
\textbf{The entire gain arrives in the final two steps, and those two are the
only steps in the ladder that improve on every seed.} This is stated plainly because it
is precisely the structure an ablation study exists to expose, and because a
reader who reconstructs it independently from Table~\ref{tab:ladder} would be
entitled to distrust anything else in this chapter.

% ===========================================================================
\section{Statistical Inference}
\label{sec:results-inference}

\subsection{A Correction to the Planned Analysis}

Section~\ref{sec:method-statistics} specifies a paired patient-clustered
bootstrap with false discovery rate control. Executed literally, that procedure
produced a table that contradicted itself: cross-sequence self-supervision
appeared as $+0.0270$ with $p = 0.000$ on seed~42 and as $-0.0234$ with
$p = 0.000$ on seed~43. Both intervals excluded zero; their signs were opposite.

The cause is that the bootstrap resamples \emph{patients} with the trained model
held fixed. Its interval therefore covers test-set sampling and excludes training
stochasticity. On this problem that omission is decisive, because the
seed-to-seed standard deviation is larger than most of the effects under test.
Reported per seed, the analysis would have permitted any claim to be supported by
selecting a seed.

Two changes were made, and neither alters what is compared.

\begin{enumerate}
  \item \emph{Seeds are averaged inside each bootstrap replicate.} One patient
  resample is drawn and applied to all seeds, the difference computed per seed,
  and the seeds averaged before the replicate is recorded. The resulting interval
  is for the mean effect over training runs. The between-seed standard deviation
  is reported alongside it rather than absorbed into it.

  \item \emph{The false discovery rate family contains one test per comparison
  and metric.} Three seeds are three views of one comparison, not three
  independent tests; including them separately tripled the apparent multiplicity.
\end{enumerate}

Per-seed intervals are retained in the appendix as descriptive quantities
carrying no significance claim.

\subsection{Primary Comparisons}

Table~\ref{tab:comparisons} reports the pre-specified comparisons under the
corrected procedure. Three survive correction, where the three-seed
analysis produced one.

\begin{table}[htbp]
\centering
\small
\caption{Pre-specified comparisons, QWK, seven seeds. Difference averaged
over seeds with the patient resample shared across them; FDR controlled
across these rows.}
\label{tab:comparisons}
\begin{tabular}{llccccc}
\toprule
Comparison & Tests & $\Delta$ & 95\% CI & Seeds $+$ & $p_{\mathrm{FDR}}$ & Sig \\
\midrule
E7 vs E0 & full system & $+0.0177$ & $[+0.0097, +0.0260]$ & 7/7 & $\mathbf{<0.001}$ & \textbf{yes} \\
E6 vs E5 & typed edges (RQ1) & $+0.0093$ & $[+0.0027, +0.0156]$ & 7/7 & $\mathbf{0.009}$ & \textbf{yes} \\
E7 vs E6 & ordinal head (RQ5) & $+0.0082$ & $[+0.0025, +0.0143]$ & 7/7 & $\mathbf{0.009}$ & \textbf{yes} \\
\midrule
E2 vs E1 & routing (RQ3) & $+0.0053$ & $[-0.0002, +0.0112]$ & 5/7 & $0.131$ & no \\
E4 vs E3 & ACSSL (RQ2) & $+0.0020$ & $[-0.0043, +0.0089]$ & 4/7 & $0.688$ & no \\
E6 vs E6\_shuf & anatomy (RQ1) & $+0.0020$ & $[-0.0038, +0.0076]$ & 4/7 & $0.688$ & no \\
E6 vs E6\_ung & gating & $+0.0014$ & $[-0.0041, +0.0064]$ & 3/7 & $0.713$ & no \\
E5 vs E0 & message passing & $+0.0002$ & $[-0.0073, +0.0074]$ & 5/7 & $0.967$ & no \\
E3 vs E2 & modality dropout & $-0.0026$ & $[-0.0091, +0.0034]$ & 4/7 & $0.627$ & no \\
\bottomrule
\end{tabular}
\end{table}

\subsection{Effect Size and Statistical Power}
\label{sec:results-power}

A p-value confounds effect size with sample size, and a null result is
uninterpretable without knowing what the study could have detected.
Table~\ref{tab:power} therefore reports Cohen's $d$ on the paired across-seed
differences alongside the number of training seeds that would be required to
detect an effect of the observed size at 80\% power.

\begin{table}[htbp]
\centering
\caption{Effect size and detectability over seven seeds. Cohen's $d$ is
computed on the paired across-seed differences. \enquote{Seeds required} is
for 80\% power at the observed effect size and between-seed variance.}
\label{tab:power}
\begin{tabular}{lcrr}
\toprule
Effect & $\Delta$ QWK & Cohen's $d$ & Seeds required \\
\midrule
Ordinal and cost-sensitive head (RQ5) & $+0.0082$ & $\mathbf{1.97}$ & 3 \\
Full system vs baseline & $+0.0177$ & $\mathbf{1.91}$ & 3 \\
Typed heterogeneous edges (RQ1) & $+0.0093$ & $\mathbf{1.06}$ & 7 \\
\midrule
Disease-conditioned routing (RQ3) & $+0.0053$ & $0.39$ & 52 \\
Anatomical topology vs shuffle (RQ1) & $+0.0020$ & $0.20$ & 205 \\
Modality dropout & $-0.0026$ & $-0.13$ & 442 \\
Cross-sequence self-supervision (RQ2) & $+0.0020$ & $0.12$ & 545 \\
Gated residual & $+0.0014$ & $0.07$ & 1{,}615 \\
Relational message passing & $+0.0002$ & $0.01$ & 65{,}489 \\
\bottomrule
\end{tabular}
\end{table}

The table separates into two groups with nothing in between. Three effects are
large by any convention, $d$ between $1.06$ and $1.97$, and all three are the
comparisons that survive correction. Everything else is $d < 0.4$. That gap
matters: a study genuinely starved of power would show a spread of intermediate
effects rather than a clean division, and would not resolve any comparison at
three seeds.

The ladder is therefore not an instrument that reports nothing everywhere. Two
effects are detectable with a single seed and a third with seven, which is what
the campaign used. Where it reports nothing, that is a statement about the
effect rather than about the instrument.

The largest standardised effect in the study is the ordinal and cost-sensitive
head at $d = 1.97$, exceeding the full-system comparison it forms part of. Its
between-seed standard deviation is $0.0042$, the smallest of any comparison, so
its modest raw difference is also the most reliably reproduced.

These figures supersede a three-seed version of the same table, and the
differences are instructive. Three-seed variance estimates are themselves noisy:
routing moved from an estimated $3{,}467$ seeds to $52$, and cross-sequence
self-supervision from $216$ to $545$. The ordering is stable and the two groups
are unchanged, but individual requirements should be read as order-of-magnitude
guidance rather than as precise targets.

\subsection{What the Nulls Establish}
\label{sec:results-equivalence}

Reporting a comparison as \enquote{not significant} states only that an effect
was not detected. It is the weakest available form of a null, and it invites the
reader to supply the explanation that the study was too small. The paired
interval already bounds the effect, and stating that bound converts an absence
into a measurement.

Table~\ref{tab:bounds} gives, for each comparison that did not separate, the far
edge of its 95\% interval: the largest effect still compatible with the data.

These are \emph{confidence} bounds and are deliberately not called equivalence
bounds. A formal equivalence test requires a smallest effect of interest to be
specified in advance and tested against, for example by two one-sided tests.
No such margin was pre-specified here, and choosing one now, with the results
already seen, would be post hoc. The bounds below therefore say what the data
are compatible with; they do not assert equivalence to zero.

\begin{table}[htbp]
\centering
\caption{Upper compatibility bounds for the comparisons that did not
separate. Each figure is the far edge of the 95\% confidence interval: the
largest effect still compatible with the data, in QWK and as a percentage
of the E0 baseline of $0.7270$. These are confidence bounds, not the
outcome of an equivalence procedure.}
\label{tab:bounds}
\begin{tabular}{lrr}
\toprule
Comparison & Effect is at most & As \% of baseline \\
\midrule
Gated residual & $0.0064$ & $0.89\%$ \\
Relational message passing & $0.0074$ & $1.02\%$ \\
Anatomical topology vs shuffle (RQ1) & $0.0076$ & $1.05\%$ \\
Cross-sequence self-supervision (RQ2) & $0.0089$ & $1.22\%$ \\
Modality dropout & $0.0091$ & $1.25\%$ \\
Disease-conditioned routing (RQ3) & $0.0112$ & $1.55\%$ \\
\bottomrule
\end{tabular}
\end{table}

The claims available are therefore considerably stronger than the absence of a
significant result. Anatomically correct graph topology contributes at most
$1.05\%$ of baseline performance over a degree-preserving random shuffle.
Anatomically aligned cross-sequence self-supervision contributes at most
$1.22\%$ over the same architecture without it. Disease-conditioned routing
contributes at most $1.55\%$ over fixed fusion.

Set against the reference standard, these bounds are small in a way that matters.
Section~\ref{sec:lr-reliability} reports inter-reader $\kappa$ between $0.49$ and
$0.73$ depending on compartment. An effect bounded below $1.6\%$ of baseline is
not merely undetectable in this study; it is below the level at which the
reference standard itself could adjudicate the question.

\subsection{The Seven-Seed Extension}
\label{sec:results-sevenseed}

The campaign was extended from three seeds to seven, and the extension is
reported rather than absorbed, because it changed two verdicts.

Typed heterogeneous edges moved from $+0.0123$ on 3/3 seeds with
$p_{\mathrm{FDR}} = 0.054$ --- narrowly missing correction --- to $+0.0093$ on
7/7 seeds with $p_{\mathrm{FDR}} = 0.009$. The point estimate fell while the
between-seed standard deviation fell further, and the comparison crossed the
threshold. The ordinal head moved the same way, from $0.054$ to $0.009$.

Anatomical topology moved in the opposite direction: from $+0.0051$ on 2/3 seeds
to $+0.0020$ on 4/7. So did cross-sequence self-supervision, from $+0.0051$ on
2/3 to $+0.0020$ on 4/7.

That divergence is the reason the extension is worth reporting. The same four
additional runs that promoted two comparisons to significance made two others
weaker. Insufficient power is therefore not available as a blanket explanation
for the nulls in this chapter: the instrument demonstrably resolved effects at
this scale, and the effects it did not resolve moved away from significance as
evidence accumulated rather than toward it.

\subsection{Clinical Error Structure}
\label{sec:results-errors}

Aggregate agreement conceals the errors that matter clinically. A
Severe$\rightarrow$Normal/Mild confusion may withhold treatment from a patient
who needs it, and the cost matrix of Section~\ref{sec:method-cost} was chosen to
suppress that direction specifically.

\begin{table}[htbp]
\centering
\caption{Clinical error structure over seven seeds. \enquote{Distance $\geq 2$}
is the rate of predictions two or more grades from the reference.}
\label{tab:errors}
\begin{tabular}{lrrr}
\toprule
Configuration & Severe recall & Severe $\rightarrow$ Normal & Distance $\geq 2$ \\
\midrule
E0 single sequence & 61.1\% & 5.7\% & 0.526\% \\
E1 multi-sequence & 61.1\% & 4.9\% & 0.500\% \\
E2 + disease routing & 60.7\% & 4.7\% & 0.414\% \\
E3 + modality dropout & 61.4\% & 4.6\% & 0.420\% \\
E4 + ACSSL & 58.8\% & 4.0\% & 0.453\% \\
E5 + homogeneous graph & 59.0\% & 4.1\% & 0.389\% \\
E6 + typed graph & 62.9\% & 4.1\% & 0.375\% \\
\quad E6 shuffled edges & 59.3\% & 3.7\% & 0.330\% \\
\quad E6 ungated residual & \textbf{65.1\%} & 4.3\% & 0.500\% \\
E7 + ordinal, cost-sensitive & 63.1\% & \textbf{3.8\%} & \textbf{0.344\%} \\
\bottomrule
\end{tabular}
\end{table}

Severe$\rightarrow$Normal errors fall from $5.7\%$ at E0 to $3.8\%$ at E7, a
relative reduction of a third, while recall on Severe targets rises from
$61.1\%$ to $63.1\%$. Distant errors fall by a comparable margin, from $0.526\%$
to $0.344\%$. The system does not purchase aggregate agreement by trading away
the errors clinicians care most about; both move in the desired direction
together.

One complication should be stated rather than smoothed over. The ungated variant
achieves the highest Severe recall of any configuration at $65.1\%$, and
simultaneously carries the joint-highest rate of distant errors at $0.500\%$
against E7's $0.344\%$. Greater sensitivity bought with more badly misplaced
predictions is precisely the trade the cost matrix exists to refuse, and E7
makes it in the intended direction.

\subsection{Calibration}
\label{sec:results-calibration}

The expected calibration error recorded for each run is the \emph{uncalibrated}
test value; temperature scaling is fitted on the validation partition and its
test metrics are stored separately. The distinction is material, because the
uncalibrated figures alone would suggest the ladder becomes progressively worse
calibrated, and the correction the protocol actually applies reverses that.

\begin{table}[htbp]
\centering
\caption{Expected calibration error before and after temperature scaling, and
the fitted temperature, over seven seeds.}
\label{tab:calibration}
\begin{tabular}{lrrr}
\toprule
Configuration & ECE uncalibrated & ECE calibrated & Temperature \\
\midrule
E0 single sequence & $0.0293$ & $\mathbf{0.0130}$ & $1.161$ \\
E1 multi-sequence & $0.0214$ & $0.0129$ & $1.131$ \\
E2 + disease routing & $0.0428$ & $0.0168$ & $1.601$ \\
E3 + modality dropout & $0.0326$ & $0.0145$ & $1.200$ \\
E4 + ACSSL & $0.0266$ & $0.0110$ & $1.187$ \\
E5 + homogeneous graph & $0.1024$ & $0.0397$ & $2.019$ \\
E6 + typed graph & $0.0426$ & $0.0209$ & $1.422$ \\
E7 + ordinal, cost-sensitive & $0.0508$ & $0.0245$ & $1.306$ \\
\bottomrule
\end{tabular}
\end{table}

Temperature scaling roughly halves ECE at every rung, and every fitted
temperature exceeds $1$, so every configuration is overconfident before
correction.

Two observations run against the system and are reported for that reason. E5,
the homogeneous graph, is by a wide margin the worst calibrated configuration at
$0.1024$ uncalibrated with a fitted temperature of $2.019$; it is also the rung
that fails to improve accuracy, and its overconfidence is a second symptom of
the same failure. And the accumulating ladder does not improve calibration: E7
at $0.0245$ is worse than E0 at $0.0130$, so the gains in agreement come with
less reliable probabilities. RQ5 asked whether calibrated uncertainty could
support selective prediction. On this evidence the agreement half of that
question is answered affirmatively and the calibration half is not.

% ===========================================================================
\section{RQ1 --- Relational Disease Modelling}
\label{sec:results-rq1}

\begin{quote}
\emph{Does representing lumbar disease targets as a typed heterogeneous graph
improve grading over independent target heads, an ordered five-level transformer
and a homogeneous graph, and are any gains concentrated where single-target
evidence is weakest?}
\end{quote}

RQ1 divides into two claims, and the evidence separates them.

\subsection{Typed Structure Helps}

A typed heterogeneous graph outperforms a homogeneous one by $+0.0093$ QWK
($[+0.0027, +0.0156]$), and does so on \textbf{every one of seven seeds}. The
comparison survives false discovery rate correction at
$p_{\mathrm{FDR}} = 0.009$, with $d = 1.06$.

This is one of the two verdicts the seven-seed extension changed. At three seeds
the same comparison stood at $+0.0123$ with $p_{\mathrm{FDR}} = 0.054$, narrowly
outside correction, and the power analysis then indicated roughly ten seeds would
be needed. Seven sufficed, because the additional runs reduced the between-seed
standard deviation faster than they reduced the point estimate.

Relational message passing alone does not help: E5 against E0 is $+0.0002$
($[-0.0073, +0.0074]$, $p_{\mathrm{FDR}} = 0.967$), which is as close to exactly
nothing as this study measures. The benefit appears only once relations are
typed. Message passing over an untyped graph is, on this evidence, worth
precisely its parameter count and no more.

\subsection{Anatomical Topology Does Not}

H3 predicted that \enquote{performance under a shuffled-edge control should
deteriorate if the topology itself matters}. It does not deteriorate
meaningfully. E6 exceeds its shuffled control by $+0.0020$ QWK
($[-0.0038, +0.0076]$, 4/7 seeds, $p_{\mathrm{FDR}} = 0.688$), and the paired
interval bounds any advantage at $1.05\%$ of baseline performance. The control
preserves node count, edge count and degree, so the two models are matched in
capacity and differ only in which targets are connected to which.

The conclusion is therefore sharper than the hypothesis anticipated.
\emph{Typing the relations carries information; the anatomical identity of the
relations does not.} What the graph supplies is a structured parameterisation of
interaction between targets, not the specific adjacency of levels, compartments
and bilateral pairs that motivated it.

\subsection{The Gating Ablation}

Chapter~\ref{ch:methodology} designates the ungated residual a mandatory
ablation. The gated and ungated formulations are indistinguishable
($+0.0014$, 3/7 seeds, $p_{\mathrm{FDR}} = 0.713$), with any difference bounded
at $0.89\%$ of baseline --- the tightest bound in the study. The gate can be
removed without cost.

\paragraph{Answer to RQ1.} Partly. Typed relational structure improves grading
over both independent heads and a homogeneous graph. The anatomical content of
that topology is not what produces the improvement, and the gating mechanism
contributes nothing.

% ===========================================================================
\section{RQ2 --- Anatomically Aligned Self-Supervision}
\label{sec:results-rq2}

\begin{quote}
\emph{Does cross-sequence pretraining based on DICOM-defined anatomical
correspondence improve grading, label efficiency and transfer?}
\end{quote}

Pretraining itself succeeded as an optimisation: validation InfoNCE reached
$1.5162$ against a chance value of $3.4657$, so the encoders demonstrably learned
to match corresponding anatomy across sequences. The question is whether that
representation is worth anything downstream. Three independent probes say it is
not.

\emph{Accuracy.} E4 against E3 is $+0.0020$ QWK ($[-0.0043, +0.0089]$, 4/7
seeds, $p_{\mathrm{FDR}} = 0.688$, $d = 0.12$), with any benefit bounded at
$1.22\%$ of baseline. The seven-seed extension moved this comparison
\emph{away} from significance, not toward it: at three seeds it stood at
$+0.0051$ on 2/3 seeds.

\emph{Cross-sequence robustness.} This is the probe on which anatomical
self-supervision should be strongest, because a representation that has learned
one anatomy across three acquisitions ought to lose less when one acquisition is
withheld. Under the controlled input ablation of
Section~\ref{sec:results-rq3-ablation}, adding ACSSL changes reliance on the
annotated sequence by $-0.0069 \pm 0.0304$ on 2/3 seeds (this probe was run on
seeds 42--44 only) --- indistinguishable
from zero. By contrast, ordinary modality dropout reduces that reliance by
$-0.0588$ on 3/3 seeds.

\emph{Attention.} Under the attribution analysis of
Section~\ref{sec:results-attribution}, E4 is the only configuration whose
attention concentration falls \emph{below} the single-sequence baseline
($-0.027$ against E0, 1/3 seeds).

\paragraph{Answer to RQ2.} No. H1 is rejected on all three axes tested. The
pretext task is learnable and is learned; the representation it produces confers
no measurable downstream benefit at this data scale. Label efficiency and
external transfer were not separately evaluated
(Section~\ref{sec:results-rq4}), so the rejection is scoped to internal grading
and robustness.

% ===========================================================================
\section{RQ3 --- Disease-Conditioned Modality Routing}
\label{sec:results-rq3}

\begin{quote}
\emph{Can one model learn target- and patient-dependent sequence weights while
remaining robust to missing and quality-degraded MRI sequences?}
\end{quote}

\subsection{The Learned Allocation}

The router learns the clinically expected allocation with complete consistency.
Across all fifteen runs in which a gate is present, foraminal targets allocate
most weight to sagittal T1, central canal targets to sagittal T2, and
subarticular targets to axial T2. The pattern reproduced 15/15 without being
encoded as a prior.

Chapter~\ref{ch:methodology} is explicit that this is insufficient:
\enquote{the gate weights are model allocations, not causal estimates of sequence
importance}.

\subsection{Controlled Input Ablation}
\label{sec:results-rq3-ablation}

The prescribed intervention withholds each sequence at test time and measures the
per-condition loss. Withholding the sequence a condition is graded from is
severely damaging --- between $0.06$ and $0.35$ QWK, an order of magnitude larger
than any difference in Table~\ref{tab:comparisons} --- while withholding the
others is near-neutral or slightly beneficial. Agreement between allocation and
intervention is 5/5 conditions on every seed.

Taken alone, this reads as strong causal support for RQ3. It is not, and the
control is what reveals why.

\emph{E1 has no router.} It fuses the three sequences with fixed weights. It
produces the same 5/5 pattern with drops as large or larger. The dependence is a
property of the dataset: each condition is \emph{annotated} on one particular
sequence, so withholding that sequence removes the only crop actually centred on
the pathology. Any model suffers, routed or not.

Isolating each step's effect on mean reliance on the annotated sequence:

The input ablation was run on seeds 42--44, before the campaign was extended,
and its figures are three-seed.

\begin{center}
\begin{tabular}{lcc}
\toprule
Step & $\Delta$ reliance & Seeds lower \\
\midrule
Add the router             & $-0.0010 \pm 0.0382$ & 1/3 \\
Add modality dropout       & $\mathbf{-0.0588 \pm 0.0561}$ & \textbf{3/3} \\
Add ACSSL                  & $-0.0069 \pm 0.0304$ & 2/3 \\
\bottomrule
\end{tabular}
\end{center}

The router changes sequence reliance by essentially zero. Its gate weights track
the annotation structure of the dataset without creating any causal dependence
beyond what fixed fusion already possesses.

\paragraph{Answer to RQ3.} The first half is yes and the second half is no. One
model does learn target-dependent sequence weights, and those weights are
clinically sensible. But they do not improve grading ($+0.0053$ QWK, 5/7 seeds,
$p_{\mathrm{FDR}} = 0.131$, bounded at $1.55\%$ of baseline)
and they do not improve robustness to missing sequences. H2 predicted that
routing would \enquote{degrade less under missing-sequence experiments than fixed
concatenation}; it does not. The graceful degradation that the system does
exhibit is attributable to modality dropout, a standard regulariser that
Chapter~\ref{ch:introduction} explicitly does not claim as original.

% ===========================================================================
\section{RQ4 --- Cross-Institutional Transfer}
\label{sec:results-rq4}

\textbf{Not yet answered.} The zero-shot transfer experiment on the Rizgary
Hospital cohort has not been executed. Reporting its status honestly is required
by the protocol changelog commitment of
Section~\ref{sec:intro-hypotheses}.

Preparatory work is complete and is documented in
\texttt{thesis/chapter4/}: the cohort has been surveyed, radiology reports parsed
into structured labels, imaging and report records reconciled, and a grading
worklist produced. Of the cases with both a report and usable imaging, sixteen
carry conflicts between the report-derived label and the spreadsheet reference
that require adjudication by a reader before any transfer number can be
computed.

Two blockers stand between the present state and an RQ4 result. First, the
sixteen conflicting cases await regrading. Second, the cohort DICOMs carry
unredacted patient identifiers, and the existing de-identification routine keys
on a \texttt{PatientID} field that is not unique in this cohort --- $45$ distinct
values across $346$ cases --- so applying it would merge roughly eight patients
into each pseudonym. Neither blocker is technical in the sense of requiring new
method; both require decisions and reader time.

H4 and H5 are therefore untested. They are retained here rather than removed, in
keeping with the requirement that components not supported by executed
experiments be reported as such rather than deleted from the research history.

% ===========================================================================
\section{RQ5 --- Clinically Asymmetric Error}
\label{sec:results-rq5}

\begin{quote}
\emph{Do ordinal and cost-sensitive objectives reduce clinically consequential
distant errors, particularly Severe$\rightarrow$Normal/Mild?}
\end{quote}

Yes, and this is the second most reliable effect in the study. Adding the ordinal
and cost-sensitive head improves QWK by $+0.0082$ ($[+0.0025, +0.0143]$,
$p_{\mathrm{FDR}} = 0.009$, $d = 1.97$ --- the largest standardised effect in the
study) on 7/7
seeds, with the smallest between-seed standard deviation of any comparison
($0.0017$) --- so small that a single seed suffices to detect it
(Table~\ref{tab:power}).

The clinically material quantity moves in the intended direction. Recall on
Severe targets rises from $62.7\%$ at E0 to $65.0\%$ at E7, and from $59.6\%$ at
E6, the same architecture under categorical cross-entropy. The objective is
doing what it was selected to do rather than trading distant errors for
aggregate accuracy.

H6 anticipated that these objectives \enquote{may reduce distant errors even if
they do not improve aggregate macro-F1}. In the event both improved.

\paragraph{Answer to RQ5.} Yes. This is noted with some irony: the objective
function, explicitly disclaimed in
Section~\ref{sec:intro-contrib-supporting} as a supporting component and not an
original method, is a larger and far more reliable contributor to final
performance than two of the three proposed novel mechanisms.

% ===========================================================================
\section{Why the Structural Priors Had Little to Add}
\label{sec:results-attribution}

The results above establish that most of the proposed mechanisms do not improve
grading. A null result with a mechanism is worth more than a null result with a
p-value, and this section supplies the mechanism.

\subsection{Method}

Each crop is centred by construction on its target's annotated coordinate, so the
pathology lies at the middle of the frame and attention on it can be measured
rather than described: the fraction of Grad-CAM mass falling inside a central
disc of $15$~mm radius at the $60$~mm field of view.

That disc covers $19.7\%$ of the frame, which is what a uniform map would score.
Because convolutional networks are centre-biased irrespective of what they have
learned, the identical architecture with untrained weights is evaluated as a
floor; it scores $0.204$.

Attribution depth was fixed on resolution grounds before results were compared. A
residual network downsamples by $32$, so at these $128$~pixel crops the final
block is $4 \times 4$ --- one cell spanning $15$~mm, too coarse to resolve the
disc. The penultimate block at $8 \times 8$ matches the $7 \times 7$ resolution
at which Grad-CAM is conventionally applied to $224$~pixel inputs, and is used
throughout. All three depths are reported in the appendix so that the dependence
is visible rather than taken on trust.

\subsection{Result}

\begin{center}
\begin{tabular}{lccc}
\toprule
Configuration & CAM mass in disc & vs.\ untrained & vs.\ E0 \\
\midrule
E0 single sequence     & $0.490 \pm 0.058$ & $+0.286$ & --- \\
E1 multi-sequence      & $0.578 \pm 0.073$ & $+0.374$ & $+0.089$ (3/3) \\
E2 + routing           & $0.580 \pm 0.063$ & $+0.376$ & $+0.090$ (2/3) \\
E3 + modality dropout  & $0.515 \pm 0.071$ & $+0.311$ & $+0.026$ (2/3) \\
E4 + ACSSL             & $0.462 \pm 0.027$ & $+0.258$ & $-0.027$ (1/3) \\
\midrule
Untrained (same architecture) & $0.204$ & --- & --- \\
Uniform map (disc area)       & $0.197$ & --- & --- \\
\bottomrule
\end{tabular}
\end{center}

Every trained configuration places between $2.3$ and $2.9$ times the chance share
of its attention on the annotated centre. Critically, \textbf{E0 already reaches
$0.490$} --- a plain single-sequence convolutional network with no routing, no
self-supervision and no graph.

This is the explanation the null results require. The structural priors of RQ1--RQ3
were motivated by the premise that anatomical context is needed to find and
interpret the lesion. The encoder locates it without them. There is little left
for a prior to contribute, and the ladder measures exactly that little.

\subsection{A Dissociation Worth Noting}

Multi-sequence input is the only step that improves attention on all three seeds
($+0.089$ against E0), and it is simultaneously the step with the most negative
accuracy effect ($-0.0055$ QWK). Additional sequences sharpen \emph{where} the
model looks without improving \emph{what it concludes}. Attention concentration
and grading accuracy are not the same quantity, and this study can separate them.

% ===========================================================================
\section{Threats to Validity}
\label{sec:results-threats}

\subsection{Errors Detected in the Analysis}

Three analyses produced incorrect conclusions before controls were added. Each is
recorded because each would have been reported as a positive finding, and because
all three erred in the direction of the hypothesis under test.

\emph{Per-seed intervals reversed sign.} Described in
Section~\ref{sec:results-inference}. Two comparisons yielded opposite,
individually significant results on different seeds.

\emph{A mechanism check without a control read as proof.} The 5/5 agreement
between routing weights and intervention (Section~\ref{sec:results-rq3-ablation})
appeared to establish RQ3 causally, with effects an order of magnitude larger
than any ladder difference, until a router-free model produced the same pattern.

\emph{An attribution probe measured the wrong encoder.} The first
implementation hooked the first encoder unconditionally, which for multi-sequence
configurations is the sagittal T1 encoder --- including on subarticular targets
graded from axial T2. It reported that no configuration improved on E0
($-0.036$, 1/3 seeds). Selecting the encoder that read each target's annotated
sequence reverses the sign to $+0.089$ on 3/3 seeds.

\subsection{Limitations}

\emph{Seed count.} Three seeds is thin for effects of this magnitude; a
seven-seed campaign is running. Table~\ref{tab:power} indicates that seven seeds
will resolve typed edges but will not resolve RQ2 or RQ3, whose requirements are
$216$ and $3{,}467$ seeds respectively.

\emph{Single benchmark.} All results are internal to LumbarDISC. RQ4 is
unanswered, so external validity is untested and no claim of generalisation is
made.

\emph{Point annotations.} Centre concentration presumes the annotated coordinate
marks the lesion. RSNA coordinates come from single readers; a systematic offset
would depress the measure uniformly across configurations.

\emph{Attribution method.} Grad-CAM is one attribution technique with known
failure modes. The untrained floor makes the comparison meaningful, but absolute
values are method-dependent.

\emph{The input ablation covers E0--E4 only.} Those rungs reach the fusion stage
directly, whereas E5--E7 route through the graph. The routing claim of RQ3 is in
any case cleanest to test at E2, with the graph out of the path.

\emph{Attribution on the graph rungs measures a different quantity for E7.}
Attribution was extended to E5--E7, which required one backward pass per node
rather than one per batch: message passing couples the nodes, so a summed
backward pass reports what an image contributed to the \emph{whole graph} rather
than to its own target, understating concentration by $0.173$ on the nodes
examined. With that corrected, E5 and E6 reach $0.341$ and $0.428$ against
untrained floors near $0.18$.

E7 pools to $0.171$, below chance, and that figure is a prevalence artefact
rather than a property of the model. Its ordinal head explains an accumulated
severity score, so asking what makes a target look severe highlights the
\emph{absence} of evidence on a target with no pathology, and $77.33\%$ of
targets are Normal/Mild. Split by grade, E7 gives $0.098$, $0.269$ and $0.415$
for Normal/Mild, Moderate and Severe, while E6 is flat at $0.429$, $0.419$ and
$0.448$. On targets that carry pathology the full system localises as well as
the rung below it. E7's pooled value is therefore not comparable with the other
rungs and is not presented as though it were.

% ===========================================================================
\section{Summary of Answers}
\label{sec:results-summary}

\begin{center}
\begin{tabular}{llp{6.6cm}}
\toprule
RQ & Verdict & Basis \\
\midrule
RQ1 & Partly & Typed edges $+0.0093$, \textbf{7/7 seeds},
$p_{\mathrm{FDR}} = 0.009$. Anatomical topology does not beat a
degree-preserving shuffle ($+0.0020$, 4/7), bounded at $1.05\%$ of baseline. \\
RQ2 & No & $+0.0020$, 4/7 seeds, bounded at $1.22\%$; no robustness benefit;
attention below baseline. Rejected on three axes, and weakened by the
seven-seed extension. \\
RQ3 & Mechanism yes, benefit no & Allocation replicates 15/15 but a router-free
model matches it under intervention. $+0.0053$ QWK, 5/7 seeds, bounded at
$1.55\%$. \\
RQ4 & Not executed & Cohort preparation complete; blocked on reader
adjudication, de-identification, and the absence of localisation coordinates. \\
RQ5 & Yes & $+0.0082$, \textbf{7/7 seeds}, $p_{\mathrm{FDR}} = 0.009$, $d = 1.97$
--- the largest standardised effect in the study. Severe$\rightarrow$Normal
errors fall $5.7\% \rightarrow 3.8\%$. \\
\midrule
System & Yes & $+0.0177$ QWK $[+0.0097, +0.0260]$, \textbf{7/7 seeds},
$p_{\mathrm{FDR}} < 0.001$. \\
\bottomrule
\end{tabular}
\end{center}

The thesis set out to determine whether anatomical structure, imposed explicitly
on a multi-sequence grading model, improves degenerative lumbar classification.
The answer this chapter supports is that it largely does not, that the
convolutional encoder already performs the localisation such structure was meant
to supply, and that the components which do improve performance are relational
typing and a clinically motivated objective rather than anatomical correspondence
or disease-conditioned routing.

Two qualifications belong with that summary. The improvement is bounded by a
reference standard on which expert readers agree at $\kappa$ between $0.49$ and
$0.73$, so its ceiling is set by the labels rather than by the architecture. And
the accumulating ladder improves agreement without improving the reliability of
its probabilities: calibrated ECE is worse at E7 than at E0.

Chapter~\ref{ch:discussion} considers what follows for the design of anatomy-aware
medical imaging models, and why mechanisms that are individually well motivated
and demonstrably learnable can nonetheless fail to pay.

\end{document}
; docmute skips this preamble.
% ===========================================================================

\documentclass[12pt,a4paper]{report}

\usepackage{lmodern}
\usepackage[T1]{fontenc}
\usepackage[utf8]{inputenc}
\usepackage[english]{babel}
\usepackage{csquotes}
\usepackage{microtype}
\usepackage[margin=2.5cm]{geometry}
\usepackage{setspace}
\onehalfspacing
\usepackage{amsmath,amssymb}
\usepackage{booktabs}
\usepackage{longtable}
\usepackage{array}
\usepackage{graphicx}
\usepackage[hidelinks]{hyperref}

\begin{document}

\chapter{Results}
\label{ch:results}

% ===========================================================================
\section{Overview}
\label{sec:results-overview}

This chapter reports the executed experimental programme. It is organised around
the research questions of Section~\ref{sec:intro-rqs} rather than around the
architecture, because the questions are what were fixed prospectively and the
architecture is only the instrument used to answer them.

Three statements summarise the chapter and are supported in the sections that
follow.

\begin{enumerate}
  \item \textbf{The complete system outperforms the single-sequence baseline.}
  Quadratic weighted kappa rises from $0.7270$ to $0.7447$, a paired difference
  of $+0.0177$ (95\% CI $[+0.0097, +0.0260]$) that holds on all seven seeds and
  survives false discovery rate correction. Two further comparisons also survive:
  typed heterogeneous edges and the ordinal, cost-sensitive head.

  \item \textbf{Most of the proposed mechanisms do not account for that gain.}
  Anatomically aligned self-supervision (RQ2) and disease-conditioned routing
  (RQ3) produce effects of $+0.0020$ and $+0.0053$ QWK, bounded by the paired
  intervals at $1.22\%$ and $1.55\%$ of baseline performance respectively.
  Relational structure (RQ1) helps, but the anatomical content of that structure
  does not: a typed graph beats a homogeneous one on every seed, while an
  anatomically correct topology does not beat a degree-preserving random shuffle
  and is bounded at $1.05\%$ of baseline.

  \item \textbf{There is a mechanism for the second statement, not merely a
  p-value.} Gradient-weighted class activation mapping shows that the
  convolutional encoder already concentrates $2.4\times$ the chance share of its
  attention on the annotated target before any structural prior is added. The
  localisation problem the priors were designed to assist is largely solved
  upstream of them.
\end{enumerate}

Section~\ref{sec:intro-conceptual} anticipated exactly this possibility in
stating that \enquote{if local evidence alone is sufficient, more complex
mechanisms should not be retained}. The results are read here in that spirit.

% ===========================================================================
\section{Experimental Execution}
\label{sec:results-execution}

\subsection{What Was Run}

The ablation ladder of Section~\ref{sec:method-ablation} was executed in full on
the LumbarDISC public benchmark: ten configurations $\times$ seven training
seeds $(42$--$48)$, fifty epochs each, seventy runs in total, completed in
$955.8$~minutes on a single RTX~5090. No run failed.

The campaign was extended from three seeds to seven after the three-seed
analysis left two comparisons close to the corrected significance threshold and
the power analysis of Section~\ref{sec:results-power} indicated roughly ten
seeds would resolve them. The extension is reported because it changed two
verdicts, and Section~\ref{sec:results-sevenseed} states what changed and in
which direction.

The ten configurations comprise the eight internal ladder rungs E0--E7 and
two controls attached to E6: \texttt{E6\_shuffled}, in which edge endpoints are permuted while
node and edge counts are preserved, and \texttt{E6\_ungated}, in which the gated
residual is replaced by an ungated one.

The ladder specified in Section~\ref{sec:method-ablation} runs to E8, the
selected complete system under external transfer. \textbf{E8 was not executed.}
It is the operational form of RQ4 and is blocked by the two obstacles reported in
Section~\ref{sec:results-rq4}; every result in this chapter is therefore internal
to LumbarDISC, and no rung of the executed ladder speaks to institutional shift.

\subsection{Data and Partitioning}

All results are computed on the frozen patient-level test partition: $297$
patients contributing $7{,}310$ scored targets. The partition is written once
and reloaded by every run, so all comparisons are paired at the target level.
Partition membership is derived from a fixed module constant and is independent
of the training seed by construction, which is verified behaviourally rather
than by inspection.

\subsection{Integrity Controls}

Three controls guard the reported numbers.

\emph{A label-shuffled negative control.} A copy of E0 trained on permuted labels
must not exceed chance. It does not, confirming that the evaluation pathway
cannot manufacture apparent skill from a broken model.

\emph{A run-configuration fingerprint.} Every result file records the
configuration that produced it. A run whose configuration has changed is
re-executed rather than reused. This was added after a stale twelve-epoch result
was silently reused in place of a fifty-epoch one.

\emph{A behavioural test suite.} One hundred and thirteen tests exercise the
components against the claims made about them in Chapter~\ref{ch:methodology},
including the ordinal head's cumulative transform, evidence masking in the graph,
split persistence, and the across-seed inference introduced in
Section~\ref{sec:results-inference}. All pass.

% ===========================================================================
\section{ROI Quality Control}
\label{sec:results-roiqc}

Section~\ref{sec:method-roi-qc} commits to visual inspection of a validation
subset using overlays of the detected level, crop boundary and corresponding
axial slices, recording correct lumbar level, inclusion of the relevant
canal/foramen/recess, correct left/right orientation, adequate axial coverage,
crop truncation, and correspondence failure caused by unusual axial angulation.
It further commits to a final inclusion/exclusion table.

This section reports that inspection. It is placed before the results because
one of its findings determines whether a later result can be interpreted at all.

\subsection{Why the Correspondence Had to Be Verified}

Core Contribution I rests entirely on DICOM-defined anatomical correspondence:
ACSSL's positive pairs are anatomically aligned rather than augmentation-derived.
If that alignment were wrong, the pretext task would have been matching
misaligned patches, and the null result for RQ2
(Section~\ref{sec:results-rq2}) would be an artefact of the data pipeline rather
than a finding about self-supervision.

No data-grounded evidence existed either way. The correspondence code records
that an earlier implementation never read \texttt{ImageOrientationPatient} and
substituted textbook direction cosines selected by a substring match on the
series description, and that its reported round-trip error of $0.00000000$~mm
was vacuous: mapping a point through a matrix and back through its inverse
returns the original point whatever the matrix contains.

\subsection{An External Validation}

The check used here compares two \emph{independent human annotations} through
the transform, so it cannot succeed by algebraic identity. LumbarDISC annotates
different conditions of the same level on different sequences: central canal on
sagittal T2, subarticular on axial T2. The sagittal canal annotation for a level
is therefore projected into the axial series and asked whether it selects the
slice that carries that level's own axial annotations.

\begin{table}[htbp]
\centering
\caption{Cross-annotation correspondence check. Offset is in whole slices from
the nearest independently annotated target at the same level.}
\label{tab:roiqc}
\begin{tabular}{lrrrrrr}
\toprule
Partition & Projections & 0 & 1 & 2 & $>2$ & Within 1 \\
\midrule
Test (297 studies)  & 1{,}429 & 69\% & 24\% & 2\% & 6\% & 93\% \\
Dev (1{,}676 studies) & 8{,}113 & 70\% & 24\% & 2\% & 5\% & \textbf{94\%} \\
All & 9{,}542 & 70\% & 24\% & 2\% & 5\% & 93.6\% \\
\bottomrule
\end{tabular}
\end{table}

Median offset is $0.0$~mm in every partition and the 90th percentile is
$4.0$~mm, one slice thickness. The dev row is the relevant one for Contribution
I, since ACSSL pretrained on that partition.

\subsection{The Residual Failures Are Not Systematic}

Whether the remaining 5--6\% undermines the verdict depends on whether failure is
concentrated and whether it is level-dependent.

Of 1{,}973 studies, 1{,}734 (88\%) have no failing projection at all, and the
nine studies that fail on every projection account for only 8\% of all outliers.
Failure is spread thin rather than confined to a broken subpopulation.

More decisively, outliers are evenly distributed across levels: 19.7\%, 20.6\%,
19.9\%, 19.1\% and 20.8\% for L1--L2 through L5--S1, against an overall
distribution of 18.8--20.6\%. A systematically incorrect affine transform would
not produce that profile; it would fail worst at L5--S1, where segmental
angulation is greatest. A flat level profile is the signature of annotation
noise and genuine acquisition variation.

The nine complete failures, with median offsets of 16--67~mm, are recorded as
the \enquote{correspondence failure caused by unusual axial angulation} category
and form the first entries of the inclusion/exclusion table.

\subsection{Left/Right Orientation}

Orientation is measurable rather than impressionistic. Across 2{,}881 axial
annotations, left subarticular targets sit at $+8.99$~mm and right subarticular
targets at $-8.14$~mm relative to each slice's own horizontal centre, where
positive is right of centre. Radiological convention renders patient-left on the
image right, so left targets appearing at $+9$~mm is correct. Separation is
$17.12$~mm and departure from a perfect mirror is $0.85$~mm.

This also constrains the lateralisation observation reported in
Section~\ref{sec:results-attribution}: since the annotations are symmetric to
within $0.85$~mm, the asymmetry observed in the attention maps is not inherited
from the data.

\subsection{Review Sheets and Reproducibility}

One sheet per study and level, five condition rows, two panels per row: the
annotated plane with a solid marker and the same three-dimensional point
projected into the orthogonal plane with a dashed marker. Markers are $10$~mm in
physical radius, sized per slice from \texttt{PixelSpacing} so they are
comparable across scanners. Colour encodes the reference grade and a red border
marks a crop truncated at the image edge.

Each panel states its \emph{native} pixel size. A $100$~mm field of view is
213~px at $0.469$~mm spacing but 107~px at $0.938$~mm, and an upsampled panel
that does not declare what it was upsampled from implies detail it does not
have.

The current set is 300 sheets over 60 studies --- 20\% of the test partition ---
at all five levels, covering 1{,}474 targets, with 0.5\% truncated crops and a
median cross-plane distance of $1.1$~mm, well inside the $3.5$--$4.8$~mm slice
thickness, so the projected point lies within the displayed slice rather than
being extrapolated to it.

Studies are drawn at random from the frozen partition under a fixed seed, so the
subset is reproducible and cannot be reselected after a result is seen. The
generating command is stored beside the figures, and a manifest records every
sheet with its study, level, target count, partition, seed and geometry
parameters.

The accompanying checklist carries one row per target with the six Chapter~3
fields left blank for a reader. Completing it produces the inclusion/exclusion
table; at the time of writing the reader pass is outstanding, and this is stated
rather than presented as complete.

% ===========================================================================
\section{The Ablation Ladder}
\label{sec:results-ladder}

\subsection{Aggregate Performance}

Table~\ref{tab:ladder} reports each configuration averaged over the seven
seeds.

\begin{table}[htbp]
\centering
\caption{Ablation ladder on the held-out test partition, mean $\pm$ standard
deviation over seven training seeds.}
\label{tab:ladder}
\begin{tabular}{lccc}
\toprule
Configuration & QWK & Macro-F1 & Balanced acc. \\
\midrule
E0 single sequence          & $0.7270 \pm 0.0108$ & $0.693 \pm 0.009$ & 69.2\% \\
E1 multi-sequence           & $0.7246 \pm 0.0075$ & $0.694 \pm 0.004$ & 68.9\% \\
E2 + disease routing        & $0.7298 \pm 0.0123$ & $0.701 \pm 0.009$ & 69.6\% \\
E3 + modality dropout       & $0.7273 \pm 0.0125$ & $0.700 \pm 0.009$ & 70.1\% \\
E4 + ACSSL                  & $0.7293 \pm 0.0109$ & $0.697 \pm 0.008$ & 69.3\% \\
E5 + homogeneous graph      & $0.7272 \pm 0.0091$ & $0.698 \pm 0.008$ & 69.1\% \\
E6 + typed graph            & $0.7364 \pm 0.0073$ & $0.706 \pm 0.010$ & 70.7\% \\
\quad E6 shuffled edges     & $0.7344 \pm 0.0052$ & $0.704 \pm 0.007$ & 69.9\% \\
\quad E6 ungated residual   & $0.7350 \pm 0.0167$ & $0.700 \pm 0.014$ & 70.6\% \\
E7 + ordinal, cost-sensitive & $\mathbf{0.7447 \pm 0.0051}$ & $\mathbf{0.712 \pm 0.005}$ & \textbf{71.0\%} \\
\bottomrule
\end{tabular}
\end{table}

\subsection{Where the Gain Actually Arrives}

Reporting only the endpoints would conceal the shape of the ladder.
Table~\ref{tab:steps} decomposes the $+0.0177$ into its seven successive
steps.

\begin{table}[htbp]
\centering
\caption{Step-by-step decomposition from E0 to E7 over seven seeds.
\enquote{Seeds $+$} counts the seeds on which the step improved QWK.}
\label{tab:steps}
\begin{tabular}{lccr}
\toprule
Step & $\Delta$ QWK & Seeds $+$ & Cumulative \\
\midrule
Multi-sequence input & $-0.0024$ & 4/7 & $-0.0024$ \\
Disease-conditioned routing & $+0.0053$ & 5/7 & $+0.0028$ \\
Modality dropout & $-0.0026$ & 4/7 & $+0.0003$ \\
Anatomical cross-sequence SSL & $+0.0020$ & 4/7 & $+0.0023$ \\
Homogeneous graph & $-0.0021$ & 4/7 & $+0.0002$ \\
\midrule
Typed heterogeneous edges & $+0.0093$ & \textbf{7/7} & $+0.0094$ \\
Ordinal and cost-sensitive head & $+0.0082$ & \textbf{7/7} & $\mathbf{+0.0177}$ \\
\bottomrule
\end{tabular}
\end{table}

Through the first five steps the accumulated system goes nowhere. It moves
within $\pm 0.003$~QWK of the baseline it began from and ends those five steps
$+0.0002$ ahead, and not one of them improves on more than five of seven seeds.
\textbf{The entire gain arrives in the final two steps, and those two are the
only steps in the ladder that improve on every seed.} This is stated plainly because it
is precisely the structure an ablation study exists to expose, and because a
reader who reconstructs it independently from Table~\ref{tab:ladder} would be
entitled to distrust anything else in this chapter.

% ===========================================================================
\section{Statistical Inference}
\label{sec:results-inference}

\subsection{A Correction to the Planned Analysis}

Section~\ref{sec:method-statistics} specifies a paired patient-clustered
bootstrap with false discovery rate control. Executed literally, that procedure
produced a table that contradicted itself: cross-sequence self-supervision
appeared as $+0.0270$ with $p = 0.000$ on seed~42 and as $-0.0234$ with
$p = 0.000$ on seed~43. Both intervals excluded zero; their signs were opposite.

The cause is that the bootstrap resamples \emph{patients} with the trained model
held fixed. Its interval therefore covers test-set sampling and excludes training
stochasticity. On this problem that omission is decisive, because the
seed-to-seed standard deviation is larger than most of the effects under test.
Reported per seed, the analysis would have permitted any claim to be supported by
selecting a seed.

Two changes were made, and neither alters what is compared.

\begin{enumerate}
  \item \emph{Seeds are averaged inside each bootstrap replicate.} One patient
  resample is drawn and applied to all seeds, the difference computed per seed,
  and the seeds averaged before the replicate is recorded. The resulting interval
  is for the mean effect over training runs. The between-seed standard deviation
  is reported alongside it rather than absorbed into it.

  \item \emph{The false discovery rate family contains one test per comparison
  and metric.} Three seeds are three views of one comparison, not three
  independent tests; including them separately tripled the apparent multiplicity.
\end{enumerate}

Per-seed intervals are retained in the appendix as descriptive quantities
carrying no significance claim.

\subsection{Primary Comparisons}

Table~\ref{tab:comparisons} reports the pre-specified comparisons under the
corrected procedure. Three survive correction, where the three-seed
analysis produced one.

\begin{table}[htbp]
\centering
\small
\caption{Pre-specified comparisons, QWK, seven seeds. Difference averaged
over seeds with the patient resample shared across them; FDR controlled
across these rows.}
\label{tab:comparisons}
\begin{tabular}{llccccc}
\toprule
Comparison & Tests & $\Delta$ & 95\% CI & Seeds $+$ & $p_{\mathrm{FDR}}$ & Sig \\
\midrule
E7 vs E0 & full system & $+0.0177$ & $[+0.0097, +0.0260]$ & 7/7 & $\mathbf{<0.001}$ & \textbf{yes} \\
E6 vs E5 & typed edges (RQ1) & $+0.0093$ & $[+0.0027, +0.0156]$ & 7/7 & $\mathbf{0.009}$ & \textbf{yes} \\
E7 vs E6 & ordinal head (RQ5) & $+0.0082$ & $[+0.0025, +0.0143]$ & 7/7 & $\mathbf{0.009}$ & \textbf{yes} \\
\midrule
E2 vs E1 & routing (RQ3) & $+0.0053$ & $[-0.0002, +0.0112]$ & 5/7 & $0.131$ & no \\
E4 vs E3 & ACSSL (RQ2) & $+0.0020$ & $[-0.0043, +0.0089]$ & 4/7 & $0.688$ & no \\
E6 vs E6\_shuf & anatomy (RQ1) & $+0.0020$ & $[-0.0038, +0.0076]$ & 4/7 & $0.688$ & no \\
E6 vs E6\_ung & gating & $+0.0014$ & $[-0.0041, +0.0064]$ & 3/7 & $0.713$ & no \\
E5 vs E0 & message passing & $+0.0002$ & $[-0.0073, +0.0074]$ & 5/7 & $0.967$ & no \\
E3 vs E2 & modality dropout & $-0.0026$ & $[-0.0091, +0.0034]$ & 4/7 & $0.627$ & no \\
\bottomrule
\end{tabular}
\end{table}

\subsection{Effect Size and Statistical Power}
\label{sec:results-power}

A p-value confounds effect size with sample size, and a null result is
uninterpretable without knowing what the study could have detected.
Table~\ref{tab:power} therefore reports Cohen's $d$ on the paired across-seed
differences alongside the number of training seeds that would be required to
detect an effect of the observed size at 80\% power.

\begin{table}[htbp]
\centering
\caption{Effect size and detectability over seven seeds. Cohen's $d$ is
computed on the paired across-seed differences. \enquote{Seeds required} is
for 80\% power at the observed effect size and between-seed variance.}
\label{tab:power}
\begin{tabular}{lcrr}
\toprule
Effect & $\Delta$ QWK & Cohen's $d$ & Seeds required \\
\midrule
Ordinal and cost-sensitive head (RQ5) & $+0.0082$ & $\mathbf{1.97}$ & 3 \\
Full system vs baseline & $+0.0177$ & $\mathbf{1.91}$ & 3 \\
Typed heterogeneous edges (RQ1) & $+0.0093$ & $\mathbf{1.06}$ & 7 \\
\midrule
Disease-conditioned routing (RQ3) & $+0.0053$ & $0.39$ & 52 \\
Anatomical topology vs shuffle (RQ1) & $+0.0020$ & $0.20$ & 205 \\
Modality dropout & $-0.0026$ & $-0.13$ & 442 \\
Cross-sequence self-supervision (RQ2) & $+0.0020$ & $0.12$ & 545 \\
Gated residual & $+0.0014$ & $0.07$ & 1{,}615 \\
Relational message passing & $+0.0002$ & $0.01$ & 65{,}489 \\
\bottomrule
\end{tabular}
\end{table}

The table separates into two groups with nothing in between. Three effects are
large by any convention, $d$ between $1.06$ and $1.97$, and all three are the
comparisons that survive correction. Everything else is $d < 0.4$. That gap
matters: a study genuinely starved of power would show a spread of intermediate
effects rather than a clean division, and would not resolve any comparison at
three seeds.

The ladder is therefore not an instrument that reports nothing everywhere. Two
effects are detectable with a single seed and a third with seven, which is what
the campaign used. Where it reports nothing, that is a statement about the
effect rather than about the instrument.

The largest standardised effect in the study is the ordinal and cost-sensitive
head at $d = 1.97$, exceeding the full-system comparison it forms part of. Its
between-seed standard deviation is $0.0042$, the smallest of any comparison, so
its modest raw difference is also the most reliably reproduced.

These figures supersede a three-seed version of the same table, and the
differences are instructive. Three-seed variance estimates are themselves noisy:
routing moved from an estimated $3{,}467$ seeds to $52$, and cross-sequence
self-supervision from $216$ to $545$. The ordering is stable and the two groups
are unchanged, but individual requirements should be read as order-of-magnitude
guidance rather than as precise targets.

\subsection{What the Nulls Establish}
\label{sec:results-equivalence}

Reporting a comparison as \enquote{not significant} states only that an effect
was not detected. It is the weakest available form of a null, and it invites the
reader to supply the explanation that the study was too small. The paired
interval already bounds the effect, and stating that bound converts an absence
into a measurement.

Table~\ref{tab:bounds} gives, for each comparison that did not separate, the far
edge of its 95\% interval: the largest effect still compatible with the data.

These are \emph{confidence} bounds and are deliberately not called equivalence
bounds. A formal equivalence test requires a smallest effect of interest to be
specified in advance and tested against, for example by two one-sided tests.
No such margin was pre-specified here, and choosing one now, with the results
already seen, would be post hoc. The bounds below therefore say what the data
are compatible with; they do not assert equivalence to zero.

\begin{table}[htbp]
\centering
\caption{Upper compatibility bounds for the comparisons that did not
separate. Each figure is the far edge of the 95\% confidence interval: the
largest effect still compatible with the data, in QWK and as a percentage
of the E0 baseline of $0.7270$. These are confidence bounds, not the
outcome of an equivalence procedure.}
\label{tab:bounds}
\begin{tabular}{lrr}
\toprule
Comparison & Effect is at most & As \% of baseline \\
\midrule
Gated residual & $0.0064$ & $0.89\%$ \\
Relational message passing & $0.0074$ & $1.02\%$ \\
Anatomical topology vs shuffle (RQ1) & $0.0076$ & $1.05\%$ \\
Cross-sequence self-supervision (RQ2) & $0.0089$ & $1.22\%$ \\
Modality dropout & $0.0091$ & $1.25\%$ \\
Disease-conditioned routing (RQ3) & $0.0112$ & $1.55\%$ \\
\bottomrule
\end{tabular}
\end{table}

The claims available are therefore considerably stronger than the absence of a
significant result. Anatomically correct graph topology contributes at most
$1.05\%$ of baseline performance over a degree-preserving random shuffle.
Anatomically aligned cross-sequence self-supervision contributes at most
$1.22\%$ over the same architecture without it. Disease-conditioned routing
contributes at most $1.55\%$ over fixed fusion.

Set against the reference standard, these bounds are small in a way that matters.
Section~\ref{sec:lr-reliability} reports inter-reader $\kappa$ between $0.49$ and
$0.73$ depending on compartment. An effect bounded below $1.6\%$ of baseline is
not merely undetectable in this study; it is below the level at which the
reference standard itself could adjudicate the question.

\subsection{The Seven-Seed Extension}
\label{sec:results-sevenseed}

The campaign was extended from three seeds to seven, and the extension is
reported rather than absorbed, because it changed two verdicts.

Typed heterogeneous edges moved from $+0.0123$ on 3/3 seeds with
$p_{\mathrm{FDR}} = 0.054$ --- narrowly missing correction --- to $+0.0093$ on
7/7 seeds with $p_{\mathrm{FDR}} = 0.009$. The point estimate fell while the
between-seed standard deviation fell further, and the comparison crossed the
threshold. The ordinal head moved the same way, from $0.054$ to $0.009$.

Anatomical topology moved in the opposite direction: from $+0.0051$ on 2/3 seeds
to $+0.0020$ on 4/7. So did cross-sequence self-supervision, from $+0.0051$ on
2/3 to $+0.0020$ on 4/7.

That divergence is the reason the extension is worth reporting. The same four
additional runs that promoted two comparisons to significance made two others
weaker. Insufficient power is therefore not available as a blanket explanation
for the nulls in this chapter: the instrument demonstrably resolved effects at
this scale, and the effects it did not resolve moved away from significance as
evidence accumulated rather than toward it.

\subsection{Clinical Error Structure}
\label{sec:results-errors}

Aggregate agreement conceals the errors that matter clinically. A
Severe$\rightarrow$Normal/Mild confusion may withhold treatment from a patient
who needs it, and the cost matrix of Section~\ref{sec:method-cost} was chosen to
suppress that direction specifically.

\begin{table}[htbp]
\centering
\caption{Clinical error structure over seven seeds. \enquote{Distance $\geq 2$}
is the rate of predictions two or more grades from the reference.}
\label{tab:errors}
\begin{tabular}{lrrr}
\toprule
Configuration & Severe recall & Severe $\rightarrow$ Normal & Distance $\geq 2$ \\
\midrule
E0 single sequence & 61.1\% & 5.7\% & 0.526\% \\
E1 multi-sequence & 61.1\% & 4.9\% & 0.500\% \\
E2 + disease routing & 60.7\% & 4.7\% & 0.414\% \\
E3 + modality dropout & 61.4\% & 4.6\% & 0.420\% \\
E4 + ACSSL & 58.8\% & 4.0\% & 0.453\% \\
E5 + homogeneous graph & 59.0\% & 4.1\% & 0.389\% \\
E6 + typed graph & 62.9\% & 4.1\% & 0.375\% \\
\quad E6 shuffled edges & 59.3\% & 3.7\% & 0.330\% \\
\quad E6 ungated residual & \textbf{65.1\%} & 4.3\% & 0.500\% \\
E7 + ordinal, cost-sensitive & 63.1\% & \textbf{3.8\%} & \textbf{0.344\%} \\
\bottomrule
\end{tabular}
\end{table}

Severe$\rightarrow$Normal errors fall from $5.7\%$ at E0 to $3.8\%$ at E7, a
relative reduction of a third, while recall on Severe targets rises from
$61.1\%$ to $63.1\%$. Distant errors fall by a comparable margin, from $0.526\%$
to $0.344\%$. The system does not purchase aggregate agreement by trading away
the errors clinicians care most about; both move in the desired direction
together.

One complication should be stated rather than smoothed over. The ungated variant
achieves the highest Severe recall of any configuration at $65.1\%$, and
simultaneously carries the joint-highest rate of distant errors at $0.500\%$
against E7's $0.344\%$. Greater sensitivity bought with more badly misplaced
predictions is precisely the trade the cost matrix exists to refuse, and E7
makes it in the intended direction.

\subsection{Calibration}
\label{sec:results-calibration}

The expected calibration error recorded for each run is the \emph{uncalibrated}
test value; temperature scaling is fitted on the validation partition and its
test metrics are stored separately. The distinction is material, because the
uncalibrated figures alone would suggest the ladder becomes progressively worse
calibrated, and the correction the protocol actually applies reverses that.

\begin{table}[htbp]
\centering
\caption{Expected calibration error before and after temperature scaling, and
the fitted temperature, over seven seeds.}
\label{tab:calibration}
\begin{tabular}{lrrr}
\toprule
Configuration & ECE uncalibrated & ECE calibrated & Temperature \\
\midrule
E0 single sequence & $0.0293$ & $\mathbf{0.0130}$ & $1.161$ \\
E1 multi-sequence & $0.0214$ & $0.0129$ & $1.131$ \\
E2 + disease routing & $0.0428$ & $0.0168$ & $1.601$ \\
E3 + modality dropout & $0.0326$ & $0.0145$ & $1.200$ \\
E4 + ACSSL & $0.0266$ & $0.0110$ & $1.187$ \\
E5 + homogeneous graph & $0.1024$ & $0.0397$ & $2.019$ \\
E6 + typed graph & $0.0426$ & $0.0209$ & $1.422$ \\
E7 + ordinal, cost-sensitive & $0.0508$ & $0.0245$ & $1.306$ \\
\bottomrule
\end{tabular}
\end{table}

Temperature scaling roughly halves ECE at every rung, and every fitted
temperature exceeds $1$, so every configuration is overconfident before
correction.

Two observations run against the system and are reported for that reason. E5,
the homogeneous graph, is by a wide margin the worst calibrated configuration at
$0.1024$ uncalibrated with a fitted temperature of $2.019$; it is also the rung
that fails to improve accuracy, and its overconfidence is a second symptom of
the same failure. And the accumulating ladder does not improve calibration: E7
at $0.0245$ is worse than E0 at $0.0130$, so the gains in agreement come with
less reliable probabilities. RQ5 asked whether calibrated uncertainty could
support selective prediction. On this evidence the agreement half of that
question is answered affirmatively and the calibration half is not.

% ===========================================================================
\section{RQ1 --- Relational Disease Modelling}
\label{sec:results-rq1}

\begin{quote}
\emph{Does representing lumbar disease targets as a typed heterogeneous graph
improve grading over independent target heads, an ordered five-level transformer
and a homogeneous graph, and are any gains concentrated where single-target
evidence is weakest?}
\end{quote}

RQ1 divides into two claims, and the evidence separates them.

\subsection{Typed Structure Helps}

A typed heterogeneous graph outperforms a homogeneous one by $+0.0093$ QWK
($[+0.0027, +0.0156]$), and does so on \textbf{every one of seven seeds}. The
comparison survives false discovery rate correction at
$p_{\mathrm{FDR}} = 0.009$, with $d = 1.06$.

This is one of the two verdicts the seven-seed extension changed. At three seeds
the same comparison stood at $+0.0123$ with $p_{\mathrm{FDR}} = 0.054$, narrowly
outside correction, and the power analysis then indicated roughly ten seeds would
be needed. Seven sufficed, because the additional runs reduced the between-seed
standard deviation faster than they reduced the point estimate.

Relational message passing alone does not help: E5 against E0 is $+0.0002$
($[-0.0073, +0.0074]$, $p_{\mathrm{FDR}} = 0.967$), which is as close to exactly
nothing as this study measures. The benefit appears only once relations are
typed. Message passing over an untyped graph is, on this evidence, worth
precisely its parameter count and no more.

\subsection{Anatomical Topology Does Not}

H3 predicted that \enquote{performance under a shuffled-edge control should
deteriorate if the topology itself matters}. It does not deteriorate
meaningfully. E6 exceeds its shuffled control by $+0.0020$ QWK
($[-0.0038, +0.0076]$, 4/7 seeds, $p_{\mathrm{FDR}} = 0.688$), and the paired
interval bounds any advantage at $1.05\%$ of baseline performance. The control
preserves node count, edge count and degree, so the two models are matched in
capacity and differ only in which targets are connected to which.

The conclusion is therefore sharper than the hypothesis anticipated.
\emph{Typing the relations carries information; the anatomical identity of the
relations does not.} What the graph supplies is a structured parameterisation of
interaction between targets, not the specific adjacency of levels, compartments
and bilateral pairs that motivated it.

\subsection{The Gating Ablation}

Chapter~\ref{ch:methodology} designates the ungated residual a mandatory
ablation. The gated and ungated formulations are indistinguishable
($+0.0014$, 3/7 seeds, $p_{\mathrm{FDR}} = 0.713$), with any difference bounded
at $0.89\%$ of baseline --- the tightest bound in the study. The gate can be
removed without cost.

\paragraph{Answer to RQ1.} Partly. Typed relational structure improves grading
over both independent heads and a homogeneous graph. The anatomical content of
that topology is not what produces the improvement, and the gating mechanism
contributes nothing.

% ===========================================================================
\section{RQ2 --- Anatomically Aligned Self-Supervision}
\label{sec:results-rq2}

\begin{quote}
\emph{Does cross-sequence pretraining based on DICOM-defined anatomical
correspondence improve grading, label efficiency and transfer?}
\end{quote}

Pretraining itself succeeded as an optimisation: validation InfoNCE reached
$1.5162$ against a chance value of $3.4657$, so the encoders demonstrably learned
to match corresponding anatomy across sequences. The question is whether that
representation is worth anything downstream. Three independent probes say it is
not.

\emph{Accuracy.} E4 against E3 is $+0.0020$ QWK ($[-0.0043, +0.0089]$, 4/7
seeds, $p_{\mathrm{FDR}} = 0.688$, $d = 0.12$), with any benefit bounded at
$1.22\%$ of baseline. The seven-seed extension moved this comparison
\emph{away} from significance, not toward it: at three seeds it stood at
$+0.0051$ on 2/3 seeds.

\emph{Cross-sequence robustness.} This is the probe on which anatomical
self-supervision should be strongest, because a representation that has learned
one anatomy across three acquisitions ought to lose less when one acquisition is
withheld. Under the controlled input ablation of
Section~\ref{sec:results-rq3-ablation}, adding ACSSL changes reliance on the
annotated sequence by $-0.0069 \pm 0.0304$ on 2/3 seeds (this probe was run on
seeds 42--44 only) --- indistinguishable
from zero. By contrast, ordinary modality dropout reduces that reliance by
$-0.0588$ on 3/3 seeds.

\emph{Attention.} Under the attribution analysis of
Section~\ref{sec:results-attribution}, E4 is the only configuration whose
attention concentration falls \emph{below} the single-sequence baseline
($-0.027$ against E0, 1/3 seeds).

\paragraph{Answer to RQ2.} No. H1 is rejected on all three axes tested. The
pretext task is learnable and is learned; the representation it produces confers
no measurable downstream benefit at this data scale. Label efficiency and
external transfer were not separately evaluated
(Section~\ref{sec:results-rq4}), so the rejection is scoped to internal grading
and robustness.

% ===========================================================================
\section{RQ3 --- Disease-Conditioned Modality Routing}
\label{sec:results-rq3}

\begin{quote}
\emph{Can one model learn target- and patient-dependent sequence weights while
remaining robust to missing and quality-degraded MRI sequences?}
\end{quote}

\subsection{The Learned Allocation}

The router learns the clinically expected allocation with complete consistency.
Across all fifteen runs in which a gate is present, foraminal targets allocate
most weight to sagittal T1, central canal targets to sagittal T2, and
subarticular targets to axial T2. The pattern reproduced 15/15 without being
encoded as a prior.

Chapter~\ref{ch:methodology} is explicit that this is insufficient:
\enquote{the gate weights are model allocations, not causal estimates of sequence
importance}.

\subsection{Controlled Input Ablation}
\label{sec:results-rq3-ablation}

The prescribed intervention withholds each sequence at test time and measures the
per-condition loss. Withholding the sequence a condition is graded from is
severely damaging --- between $0.06$ and $0.35$ QWK, an order of magnitude larger
than any difference in Table~\ref{tab:comparisons} --- while withholding the
others is near-neutral or slightly beneficial. Agreement between allocation and
intervention is 5/5 conditions on every seed.

Taken alone, this reads as strong causal support for RQ3. It is not, and the
control is what reveals why.

\emph{E1 has no router.} It fuses the three sequences with fixed weights. It
produces the same 5/5 pattern with drops as large or larger. The dependence is a
property of the dataset: each condition is \emph{annotated} on one particular
sequence, so withholding that sequence removes the only crop actually centred on
the pathology. Any model suffers, routed or not.

Isolating each step's effect on mean reliance on the annotated sequence:

The input ablation was run on seeds 42--44, before the campaign was extended,
and its figures are three-seed.

\begin{center}
\begin{tabular}{lcc}
\toprule
Step & $\Delta$ reliance & Seeds lower \\
\midrule
Add the router             & $-0.0010 \pm 0.0382$ & 1/3 \\
Add modality dropout       & $\mathbf{-0.0588 \pm 0.0561}$ & \textbf{3/3} \\
Add ACSSL                  & $-0.0069 \pm 0.0304$ & 2/3 \\
\bottomrule
\end{tabular}
\end{center}

The router changes sequence reliance by essentially zero. Its gate weights track
the annotation structure of the dataset without creating any causal dependence
beyond what fixed fusion already possesses.

\paragraph{Answer to RQ3.} The first half is yes and the second half is no. One
model does learn target-dependent sequence weights, and those weights are
clinically sensible. But they do not improve grading ($+0.0053$ QWK, 5/7 seeds,
$p_{\mathrm{FDR}} = 0.131$, bounded at $1.55\%$ of baseline)
and they do not improve robustness to missing sequences. H2 predicted that
routing would \enquote{degrade less under missing-sequence experiments than fixed
concatenation}; it does not. The graceful degradation that the system does
exhibit is attributable to modality dropout, a standard regulariser that
Chapter~\ref{ch:introduction} explicitly does not claim as original.

% ===========================================================================
\section{RQ4 --- Cross-Institutional Transfer}
\label{sec:results-rq4}

\textbf{Not yet answered.} The zero-shot transfer experiment on the Rizgary
Hospital cohort has not been executed. Reporting its status honestly is required
by the protocol changelog commitment of
Section~\ref{sec:intro-hypotheses}.

Preparatory work is complete and is documented in
\texttt{thesis/chapter4/}: the cohort has been surveyed, radiology reports parsed
into structured labels, imaging and report records reconciled, and a grading
worklist produced. Of the cases with both a report and usable imaging, sixteen
carry conflicts between the report-derived label and the spreadsheet reference
that require adjudication by a reader before any transfer number can be
computed.

Two blockers stand between the present state and an RQ4 result. First, the
sixteen conflicting cases await regrading. Second, the cohort DICOMs carry
unredacted patient identifiers, and the existing de-identification routine keys
on a \texttt{PatientID} field that is not unique in this cohort --- $45$ distinct
values across $346$ cases --- so applying it would merge roughly eight patients
into each pseudonym. Neither blocker is technical in the sense of requiring new
method; both require decisions and reader time.

H4 and H5 are therefore untested. They are retained here rather than removed, in
keeping with the requirement that components not supported by executed
experiments be reported as such rather than deleted from the research history.

% ===========================================================================
\section{RQ5 --- Clinically Asymmetric Error}
\label{sec:results-rq5}

\begin{quote}
\emph{Do ordinal and cost-sensitive objectives reduce clinically consequential
distant errors, particularly Severe$\rightarrow$Normal/Mild?}
\end{quote}

Yes, and this is the second most reliable effect in the study. Adding the ordinal
and cost-sensitive head improves QWK by $+0.0082$ ($[+0.0025, +0.0143]$,
$p_{\mathrm{FDR}} = 0.009$, $d = 1.97$ --- the largest standardised effect in the
study) on 7/7
seeds, with the smallest between-seed standard deviation of any comparison
($0.0017$) --- so small that a single seed suffices to detect it
(Table~\ref{tab:power}).

The clinically material quantity moves in the intended direction. Recall on
Severe targets rises from $62.7\%$ at E0 to $65.0\%$ at E7, and from $59.6\%$ at
E6, the same architecture under categorical cross-entropy. The objective is
doing what it was selected to do rather than trading distant errors for
aggregate accuracy.

H6 anticipated that these objectives \enquote{may reduce distant errors even if
they do not improve aggregate macro-F1}. In the event both improved.

\paragraph{Answer to RQ5.} Yes. This is noted with some irony: the objective
function, explicitly disclaimed in
Section~\ref{sec:intro-contrib-supporting} as a supporting component and not an
original method, is a larger and far more reliable contributor to final
performance than two of the three proposed novel mechanisms.

% ===========================================================================
\section{Why the Structural Priors Had Little to Add}
\label{sec:results-attribution}

The results above establish that most of the proposed mechanisms do not improve
grading. A null result with a mechanism is worth more than a null result with a
p-value, and this section supplies the mechanism.

\subsection{Method}

Each crop is centred by construction on its target's annotated coordinate, so the
pathology lies at the middle of the frame and attention on it can be measured
rather than described: the fraction of Grad-CAM mass falling inside a central
disc of $15$~mm radius at the $60$~mm field of view.

That disc covers $19.7\%$ of the frame, which is what a uniform map would score.
Because convolutional networks are centre-biased irrespective of what they have
learned, the identical architecture with untrained weights is evaluated as a
floor; it scores $0.204$.

Attribution depth was fixed on resolution grounds before results were compared. A
residual network downsamples by $32$, so at these $128$~pixel crops the final
block is $4 \times 4$ --- one cell spanning $15$~mm, too coarse to resolve the
disc. The penultimate block at $8 \times 8$ matches the $7 \times 7$ resolution
at which Grad-CAM is conventionally applied to $224$~pixel inputs, and is used
throughout. All three depths are reported in the appendix so that the dependence
is visible rather than taken on trust.

\subsection{Result}

\begin{center}
\begin{tabular}{lccc}
\toprule
Configuration & CAM mass in disc & vs.\ untrained & vs.\ E0 \\
\midrule
E0 single sequence     & $0.490 \pm 0.058$ & $+0.286$ & --- \\
E1 multi-sequence      & $0.578 \pm 0.073$ & $+0.374$ & $+0.089$ (3/3) \\
E2 + routing           & $0.580 \pm 0.063$ & $+0.376$ & $+0.090$ (2/3) \\
E3 + modality dropout  & $0.515 \pm 0.071$ & $+0.311$ & $+0.026$ (2/3) \\
E4 + ACSSL             & $0.462 \pm 0.027$ & $+0.258$ & $-0.027$ (1/3) \\
\midrule
Untrained (same architecture) & $0.204$ & --- & --- \\
Uniform map (disc area)       & $0.197$ & --- & --- \\
\bottomrule
\end{tabular}
\end{center}

Every trained configuration places between $2.3$ and $2.9$ times the chance share
of its attention on the annotated centre. Critically, \textbf{E0 already reaches
$0.490$} --- a plain single-sequence convolutional network with no routing, no
self-supervision and no graph.

This is the explanation the null results require. The structural priors of RQ1--RQ3
were motivated by the premise that anatomical context is needed to find and
interpret the lesion. The encoder locates it without them. There is little left
for a prior to contribute, and the ladder measures exactly that little.

\subsection{A Dissociation Worth Noting}

Multi-sequence input is the only step that improves attention on all three seeds
($+0.089$ against E0), and it is simultaneously the step with the most negative
accuracy effect ($-0.0055$ QWK). Additional sequences sharpen \emph{where} the
model looks without improving \emph{what it concludes}. Attention concentration
and grading accuracy are not the same quantity, and this study can separate them.

% ===========================================================================
\section{Threats to Validity}
\label{sec:results-threats}

\subsection{Errors Detected in the Analysis}

Three analyses produced incorrect conclusions before controls were added. Each is
recorded because each would have been reported as a positive finding, and because
all three erred in the direction of the hypothesis under test.

\emph{Per-seed intervals reversed sign.} Described in
Section~\ref{sec:results-inference}. Two comparisons yielded opposite,
individually significant results on different seeds.

\emph{A mechanism check without a control read as proof.} The 5/5 agreement
between routing weights and intervention (Section~\ref{sec:results-rq3-ablation})
appeared to establish RQ3 causally, with effects an order of magnitude larger
than any ladder difference, until a router-free model produced the same pattern.

\emph{An attribution probe measured the wrong encoder.} The first
implementation hooked the first encoder unconditionally, which for multi-sequence
configurations is the sagittal T1 encoder --- including on subarticular targets
graded from axial T2. It reported that no configuration improved on E0
($-0.036$, 1/3 seeds). Selecting the encoder that read each target's annotated
sequence reverses the sign to $+0.089$ on 3/3 seeds.

\subsection{Limitations}

\emph{Seed count.} Three seeds is thin for effects of this magnitude; a
seven-seed campaign is running. Table~\ref{tab:power} indicates that seven seeds
will resolve typed edges but will not resolve RQ2 or RQ3, whose requirements are
$216$ and $3{,}467$ seeds respectively.

\emph{Single benchmark.} All results are internal to LumbarDISC. RQ4 is
unanswered, so external validity is untested and no claim of generalisation is
made.

\emph{Point annotations.} Centre concentration presumes the annotated coordinate
marks the lesion. RSNA coordinates come from single readers; a systematic offset
would depress the measure uniformly across configurations.

\emph{Attribution method.} Grad-CAM is one attribution technique with known
failure modes. The untrained floor makes the comparison meaningful, but absolute
values are method-dependent.

\emph{The input ablation covers E0--E4 only.} Those rungs reach the fusion stage
directly, whereas E5--E7 route through the graph. The routing claim of RQ3 is in
any case cleanest to test at E2, with the graph out of the path.

\emph{Attribution on the graph rungs measures a different quantity for E7.}
Attribution was extended to E5--E7, which required one backward pass per node
rather than one per batch: message passing couples the nodes, so a summed
backward pass reports what an image contributed to the \emph{whole graph} rather
than to its own target, understating concentration by $0.173$ on the nodes
examined. With that corrected, E5 and E6 reach $0.341$ and $0.428$ against
untrained floors near $0.18$.

E7 pools to $0.171$, below chance, and that figure is a prevalence artefact
rather than a property of the model. Its ordinal head explains an accumulated
severity score, so asking what makes a target look severe highlights the
\emph{absence} of evidence on a target with no pathology, and $77.33\%$ of
targets are Normal/Mild. Split by grade, E7 gives $0.098$, $0.269$ and $0.415$
for Normal/Mild, Moderate and Severe, while E6 is flat at $0.429$, $0.419$ and
$0.448$. On targets that carry pathology the full system localises as well as
the rung below it. E7's pooled value is therefore not comparable with the other
rungs and is not presented as though it were.

% ===========================================================================
\section{Summary of Answers}
\label{sec:results-summary}

\begin{center}
\begin{tabular}{llp{6.6cm}}
\toprule
RQ & Verdict & Basis \\
\midrule
RQ1 & Partly & Typed edges $+0.0093$, \textbf{7/7 seeds},
$p_{\mathrm{FDR}} = 0.009$. Anatomical topology does not beat a
degree-preserving shuffle ($+0.0020$, 4/7), bounded at $1.05\%$ of baseline. \\
RQ2 & No & $+0.0020$, 4/7 seeds, bounded at $1.22\%$; no robustness benefit;
attention below baseline. Rejected on three axes, and weakened by the
seven-seed extension. \\
RQ3 & Mechanism yes, benefit no & Allocation replicates 15/15 but a router-free
model matches it under intervention. $+0.0053$ QWK, 5/7 seeds, bounded at
$1.55\%$. \\
RQ4 & Not executed & Cohort preparation complete; blocked on reader
adjudication, de-identification, and the absence of localisation coordinates. \\
RQ5 & Yes & $+0.0082$, \textbf{7/7 seeds}, $p_{\mathrm{FDR}} = 0.009$, $d = 1.97$
--- the largest standardised effect in the study. Severe$\rightarrow$Normal
errors fall $5.7\% \rightarrow 3.8\%$. \\
\midrule
System & Yes & $+0.0177$ QWK $[+0.0097, +0.0260]$, \textbf{7/7 seeds},
$p_{\mathrm{FDR}} < 0.001$. \\
\bottomrule
\end{tabular}
\end{center}

The thesis set out to determine whether anatomical structure, imposed explicitly
on a multi-sequence grading model, improves degenerative lumbar classification.
The answer this chapter supports is that it largely does not, that the
convolutional encoder already performs the localisation such structure was meant
to supply, and that the components which do improve performance are relational
typing and a clinically motivated objective rather than anatomical correspondence
or disease-conditioned routing.

Two qualifications belong with that summary. The improvement is bounded by a
reference standard on which expert readers agree at $\kappa$ between $0.49$ and
$0.73$, so its ceiling is set by the labels rather than by the architecture. And
the accumulating ladder improves agreement without improving the reliability of
its probabilities: calibrated ECE is worse at E7 than at E0.

Chapter~\ref{ch:discussion} considers what follows for the design of anatomy-aware
medical imaging models, and why mechanisms that are individually well motivated
and demonstrably learnable can nonetheless fail to pay.

\end{document}
; docmute skips this preamble.
% ===========================================================================

\documentclass[12pt,a4paper]{report}

\usepackage{lmodern}
\usepackage[T1]{fontenc}
\usepackage[utf8]{inputenc}
\usepackage[english]{babel}
\usepackage{csquotes}
\usepackage{microtype}
\usepackage[margin=2.5cm]{geometry}
\usepackage{setspace}
\onehalfspacing
\usepackage{amsmath,amssymb}
\usepackage{booktabs}
\usepackage{longtable}
\usepackage{array}
\usepackage{graphicx}
\usepackage[hidelinks]{hyperref}

\begin{document}

\chapter{Results}
\label{ch:results}

% ===========================================================================
\section{Overview}
\label{sec:results-overview}

This chapter reports the executed experimental programme. It is organised around
the research questions of Section~\ref{sec:intro-rqs} rather than around the
architecture, because the questions are what were fixed prospectively and the
architecture is only the instrument used to answer them.

Three statements summarise the chapter and are supported in the sections that
follow.

\begin{enumerate}
  \item \textbf{The complete system outperforms the single-sequence baseline.}
  Quadratic weighted kappa rises from $0.7270$ to $0.7447$, a paired difference
  of $+0.0177$ (95\% CI $[+0.0097, +0.0260]$) that holds on all seven seeds and
  survives false discovery rate correction. Two further comparisons also survive:
  typed heterogeneous edges and the ordinal, cost-sensitive head.

  \item \textbf{Most of the proposed mechanisms do not account for that gain.}
  Anatomically aligned self-supervision (RQ2) and disease-conditioned routing
  (RQ3) produce effects of $+0.0020$ and $+0.0053$ QWK, bounded by the paired
  intervals at $1.22\%$ and $1.55\%$ of baseline performance respectively.
  Relational structure (RQ1) helps, but the anatomical content of that structure
  does not: a typed graph beats a homogeneous one on every seed, while an
  anatomically correct topology does not beat a degree-preserving random shuffle
  and is bounded at $1.05\%$ of baseline.

  \item \textbf{There is a mechanism for the second statement, not merely a
  p-value.} Gradient-weighted class activation mapping shows that the
  convolutional encoder already concentrates $2.4\times$ the chance share of its
  attention on the annotated target before any structural prior is added. The
  localisation problem the priors were designed to assist is largely solved
  upstream of them.
\end{enumerate}

Section~\ref{sec:intro-conceptual} anticipated exactly this possibility in
stating that \enquote{if local evidence alone is sufficient, more complex
mechanisms should not be retained}. The results are read here in that spirit.

% ===========================================================================
\section{Experimental Execution}
\label{sec:results-execution}

\subsection{What Was Run}

The ablation ladder of Section~\ref{sec:method-ablation} was executed in full on
the LumbarDISC public benchmark: ten configurations $\times$ seven training
seeds $(42$--$48)$, fifty epochs each, seventy runs in total, completed in
$955.8$~minutes on a single RTX~5090. No run failed.

The campaign was extended from three seeds to seven after the three-seed
analysis left two comparisons close to the corrected significance threshold and
the power analysis of Section~\ref{sec:results-power} indicated roughly ten
seeds would resolve them. The extension is reported because it changed two
verdicts, and Section~\ref{sec:results-sevenseed} states what changed and in
which direction.

The ten configurations comprise the eight internal ladder rungs E0--E7 and
two controls attached to E6: \texttt{E6\_shuffled}, in which edge endpoints are permuted while
node and edge counts are preserved, and \texttt{E6\_ungated}, in which the gated
residual is replaced by an ungated one.

The ladder specified in Section~\ref{sec:method-ablation} runs to E8, the
selected complete system under external transfer. \textbf{E8 was not executed.}
It is the operational form of RQ4 and is blocked by the two obstacles reported in
Section~\ref{sec:results-rq4}; every result in this chapter is therefore internal
to LumbarDISC, and no rung of the executed ladder speaks to institutional shift.

\subsection{Data and Partitioning}

All results are computed on the frozen patient-level test partition: $297$
patients contributing $7{,}310$ scored targets. The partition is written once
and reloaded by every run, so all comparisons are paired at the target level.
Partition membership is derived from a fixed module constant and is independent
of the training seed by construction, which is verified behaviourally rather
than by inspection.

\subsection{Integrity Controls}

Three controls guard the reported numbers.

\emph{A label-shuffled negative control.} A copy of E0 trained on permuted labels
must not exceed chance. It does not, confirming that the evaluation pathway
cannot manufacture apparent skill from a broken model.

\emph{A run-configuration fingerprint.} Every result file records the
configuration that produced it. A run whose configuration has changed is
re-executed rather than reused. This was added after a stale twelve-epoch result
was silently reused in place of a fifty-epoch one.

\emph{A behavioural test suite.} One hundred and thirteen tests exercise the
components against the claims made about them in Chapter~\ref{ch:methodology},
including the ordinal head's cumulative transform, evidence masking in the graph,
split persistence, and the across-seed inference introduced in
Section~\ref{sec:results-inference}. All pass.

% ===========================================================================
\section{ROI Quality Control}
\label{sec:results-roiqc}

Section~\ref{sec:method-roi-qc} commits to visual inspection of a validation
subset using overlays of the detected level, crop boundary and corresponding
axial slices, recording correct lumbar level, inclusion of the relevant
canal/foramen/recess, correct left/right orientation, adequate axial coverage,
crop truncation, and correspondence failure caused by unusual axial angulation.
It further commits to a final inclusion/exclusion table.

This section reports that inspection. It is placed before the results because
one of its findings determines whether a later result can be interpreted at all.

\subsection{Why the Correspondence Had to Be Verified}

Core Contribution I rests entirely on DICOM-defined anatomical correspondence:
ACSSL's positive pairs are anatomically aligned rather than augmentation-derived.
If that alignment were wrong, the pretext task would have been matching
misaligned patches, and the null result for RQ2
(Section~\ref{sec:results-rq2}) would be an artefact of the data pipeline rather
than a finding about self-supervision.

No data-grounded evidence existed either way. The correspondence code records
that an earlier implementation never read \texttt{ImageOrientationPatient} and
substituted textbook direction cosines selected by a substring match on the
series description, and that its reported round-trip error of $0.00000000$~mm
was vacuous: mapping a point through a matrix and back through its inverse
returns the original point whatever the matrix contains.

\subsection{An External Validation}

The check used here compares two \emph{independent human annotations} through
the transform, so it cannot succeed by algebraic identity. LumbarDISC annotates
different conditions of the same level on different sequences: central canal on
sagittal T2, subarticular on axial T2. The sagittal canal annotation for a level
is therefore projected into the axial series and asked whether it selects the
slice that carries that level's own axial annotations.

\begin{table}[htbp]
\centering
\caption{Cross-annotation correspondence check. Offset is in whole slices from
the nearest independently annotated target at the same level.}
\label{tab:roiqc}
\begin{tabular}{lrrrrrr}
\toprule
Partition & Projections & 0 & 1 & 2 & $>2$ & Within 1 \\
\midrule
Test (297 studies)  & 1{,}429 & 69\% & 24\% & 2\% & 6\% & 93\% \\
Dev (1{,}676 studies) & 8{,}113 & 70\% & 24\% & 2\% & 5\% & \textbf{94\%} \\
All & 9{,}542 & 70\% & 24\% & 2\% & 5\% & 93.6\% \\
\bottomrule
\end{tabular}
\end{table}

Median offset is $0.0$~mm in every partition and the 90th percentile is
$4.0$~mm, one slice thickness. The dev row is the relevant one for Contribution
I, since ACSSL pretrained on that partition.

\subsection{The Residual Failures Are Not Systematic}

Whether the remaining 5--6\% undermines the verdict depends on whether failure is
concentrated and whether it is level-dependent.

Of 1{,}973 studies, 1{,}734 (88\%) have no failing projection at all, and the
nine studies that fail on every projection account for only 8\% of all outliers.
Failure is spread thin rather than confined to a broken subpopulation.

More decisively, outliers are evenly distributed across levels: 19.7\%, 20.6\%,
19.9\%, 19.1\% and 20.8\% for L1--L2 through L5--S1, against an overall
distribution of 18.8--20.6\%. A systematically incorrect affine transform would
not produce that profile; it would fail worst at L5--S1, where segmental
angulation is greatest. A flat level profile is the signature of annotation
noise and genuine acquisition variation.

The nine complete failures, with median offsets of 16--67~mm, are recorded as
the \enquote{correspondence failure caused by unusual axial angulation} category
and form the first entries of the inclusion/exclusion table.

\subsection{Left/Right Orientation}

Orientation is measurable rather than impressionistic. Across 2{,}881 axial
annotations, left subarticular targets sit at $+8.99$~mm and right subarticular
targets at $-8.14$~mm relative to each slice's own horizontal centre, where
positive is right of centre. Radiological convention renders patient-left on the
image right, so left targets appearing at $+9$~mm is correct. Separation is
$17.12$~mm and departure from a perfect mirror is $0.85$~mm.

This also constrains the lateralisation observation reported in
Section~\ref{sec:results-attribution}: since the annotations are symmetric to
within $0.85$~mm, the asymmetry observed in the attention maps is not inherited
from the data.

\subsection{Review Sheets and Reproducibility}

One sheet per study and level, five condition rows, two panels per row: the
annotated plane with a solid marker and the same three-dimensional point
projected into the orthogonal plane with a dashed marker. Markers are $10$~mm in
physical radius, sized per slice from \texttt{PixelSpacing} so they are
comparable across scanners. Colour encodes the reference grade and a red border
marks a crop truncated at the image edge.

Each panel states its \emph{native} pixel size. A $100$~mm field of view is
213~px at $0.469$~mm spacing but 107~px at $0.938$~mm, and an upsampled panel
that does not declare what it was upsampled from implies detail it does not
have.

The current set is 300 sheets over 60 studies --- 20\% of the test partition ---
at all five levels, covering 1{,}474 targets, with 0.5\% truncated crops and a
median cross-plane distance of $1.1$~mm, well inside the $3.5$--$4.8$~mm slice
thickness, so the projected point lies within the displayed slice rather than
being extrapolated to it.

Studies are drawn at random from the frozen partition under a fixed seed, so the
subset is reproducible and cannot be reselected after a result is seen. The
generating command is stored beside the figures, and a manifest records every
sheet with its study, level, target count, partition, seed and geometry
parameters.

The accompanying checklist carries one row per target with the six Chapter~3
fields left blank for a reader. Completing it produces the inclusion/exclusion
table; at the time of writing the reader pass is outstanding, and this is stated
rather than presented as complete.

% ===========================================================================
\section{The Ablation Ladder}
\label{sec:results-ladder}

\subsection{Aggregate Performance}

Table~\ref{tab:ladder} reports each configuration averaged over the seven
seeds.

\begin{table}[htbp]
\centering
\caption{Ablation ladder on the held-out test partition, mean $\pm$ standard
deviation over seven training seeds.}
\label{tab:ladder}
\begin{tabular}{lccc}
\toprule
Configuration & QWK & Macro-F1 & Balanced acc. \\
\midrule
E0 single sequence          & $0.7270 \pm 0.0108$ & $0.693 \pm 0.009$ & 69.2\% \\
E1 multi-sequence           & $0.7246 \pm 0.0075$ & $0.694 \pm 0.004$ & 68.9\% \\
E2 + disease routing        & $0.7298 \pm 0.0123$ & $0.701 \pm 0.009$ & 69.6\% \\
E3 + modality dropout       & $0.7273 \pm 0.0125$ & $0.700 \pm 0.009$ & 70.1\% \\
E4 + ACSSL                  & $0.7293 \pm 0.0109$ & $0.697 \pm 0.008$ & 69.3\% \\
E5 + homogeneous graph      & $0.7272 \pm 0.0091$ & $0.698 \pm 0.008$ & 69.1\% \\
E6 + typed graph            & $0.7364 \pm 0.0073$ & $0.706 \pm 0.010$ & 70.7\% \\
\quad E6 shuffled edges     & $0.7344 \pm 0.0052$ & $0.704 \pm 0.007$ & 69.9\% \\
\quad E6 ungated residual   & $0.7350 \pm 0.0167$ & $0.700 \pm 0.014$ & 70.6\% \\
E7 + ordinal, cost-sensitive & $\mathbf{0.7447 \pm 0.0051}$ & $\mathbf{0.712 \pm 0.005}$ & \textbf{71.0\%} \\
\bottomrule
\end{tabular}
\end{table}

\subsection{Where the Gain Actually Arrives}

Reporting only the endpoints would conceal the shape of the ladder.
Table~\ref{tab:steps} decomposes the $+0.0177$ into its seven successive
steps.

\begin{table}[htbp]
\centering
\caption{Step-by-step decomposition from E0 to E7 over seven seeds.
\enquote{Seeds $+$} counts the seeds on which the step improved QWK.}
\label{tab:steps}
\begin{tabular}{lccr}
\toprule
Step & $\Delta$ QWK & Seeds $+$ & Cumulative \\
\midrule
Multi-sequence input & $-0.0024$ & 4/7 & $-0.0024$ \\
Disease-conditioned routing & $+0.0053$ & 5/7 & $+0.0028$ \\
Modality dropout & $-0.0026$ & 4/7 & $+0.0003$ \\
Anatomical cross-sequence SSL & $+0.0020$ & 4/7 & $+0.0023$ \\
Homogeneous graph & $-0.0021$ & 4/7 & $+0.0002$ \\
\midrule
Typed heterogeneous edges & $+0.0093$ & \textbf{7/7} & $+0.0094$ \\
Ordinal and cost-sensitive head & $+0.0082$ & \textbf{7/7} & $\mathbf{+0.0177}$ \\
\bottomrule
\end{tabular}
\end{table}

Through the first five steps the accumulated system goes nowhere. It moves
within $\pm 0.003$~QWK of the baseline it began from and ends those five steps
$+0.0002$ ahead, and not one of them improves on more than five of seven seeds.
\textbf{The entire gain arrives in the final two steps, and those two are the
only steps in the ladder that improve on every seed.} This is stated plainly because it
is precisely the structure an ablation study exists to expose, and because a
reader who reconstructs it independently from Table~\ref{tab:ladder} would be
entitled to distrust anything else in this chapter.

% ===========================================================================
\section{Statistical Inference}
\label{sec:results-inference}

\subsection{A Correction to the Planned Analysis}

Section~\ref{sec:method-statistics} specifies a paired patient-clustered
bootstrap with false discovery rate control. Executed literally, that procedure
produced a table that contradicted itself: cross-sequence self-supervision
appeared as $+0.0270$ with $p = 0.000$ on seed~42 and as $-0.0234$ with
$p = 0.000$ on seed~43. Both intervals excluded zero; their signs were opposite.

The cause is that the bootstrap resamples \emph{patients} with the trained model
held fixed. Its interval therefore covers test-set sampling and excludes training
stochasticity. On this problem that omission is decisive, because the
seed-to-seed standard deviation is larger than most of the effects under test.
Reported per seed, the analysis would have permitted any claim to be supported by
selecting a seed.

Two changes were made, and neither alters what is compared.

\begin{enumerate}
  \item \emph{Seeds are averaged inside each bootstrap replicate.} One patient
  resample is drawn and applied to all seeds, the difference computed per seed,
  and the seeds averaged before the replicate is recorded. The resulting interval
  is for the mean effect over training runs. The between-seed standard deviation
  is reported alongside it rather than absorbed into it.

  \item \emph{The false discovery rate family contains one test per comparison
  and metric.} Three seeds are three views of one comparison, not three
  independent tests; including them separately tripled the apparent multiplicity.
\end{enumerate}

Per-seed intervals are retained in the appendix as descriptive quantities
carrying no significance claim.

\subsection{Primary Comparisons}

Table~\ref{tab:comparisons} reports the pre-specified comparisons under the
corrected procedure. Three survive correction, where the three-seed
analysis produced one.

\begin{table}[htbp]
\centering
\small
\caption{Pre-specified comparisons, QWK, seven seeds. Difference averaged
over seeds with the patient resample shared across them; FDR controlled
across these rows.}
\label{tab:comparisons}
\begin{tabular}{llccccc}
\toprule
Comparison & Tests & $\Delta$ & 95\% CI & Seeds $+$ & $p_{\mathrm{FDR}}$ & Sig \\
\midrule
E7 vs E0 & full system & $+0.0177$ & $[+0.0097, +0.0260]$ & 7/7 & $\mathbf{<0.001}$ & \textbf{yes} \\
E6 vs E5 & typed edges (RQ1) & $+0.0093$ & $[+0.0027, +0.0156]$ & 7/7 & $\mathbf{0.009}$ & \textbf{yes} \\
E7 vs E6 & ordinal head (RQ5) & $+0.0082$ & $[+0.0025, +0.0143]$ & 7/7 & $\mathbf{0.009}$ & \textbf{yes} \\
\midrule
E2 vs E1 & routing (RQ3) & $+0.0053$ & $[-0.0002, +0.0112]$ & 5/7 & $0.131$ & no \\
E4 vs E3 & ACSSL (RQ2) & $+0.0020$ & $[-0.0043, +0.0089]$ & 4/7 & $0.688$ & no \\
E6 vs E6\_shuf & anatomy (RQ1) & $+0.0020$ & $[-0.0038, +0.0076]$ & 4/7 & $0.688$ & no \\
E6 vs E6\_ung & gating & $+0.0014$ & $[-0.0041, +0.0064]$ & 3/7 & $0.713$ & no \\
E5 vs E0 & message passing & $+0.0002$ & $[-0.0073, +0.0074]$ & 5/7 & $0.967$ & no \\
E3 vs E2 & modality dropout & $-0.0026$ & $[-0.0091, +0.0034]$ & 4/7 & $0.627$ & no \\
\bottomrule
\end{tabular}
\end{table}

\subsection{Effect Size and Statistical Power}
\label{sec:results-power}

A p-value confounds effect size with sample size, and a null result is
uninterpretable without knowing what the study could have detected.
Table~\ref{tab:power} therefore reports Cohen's $d$ on the paired across-seed
differences alongside the number of training seeds that would be required to
detect an effect of the observed size at 80\% power.

\begin{table}[htbp]
\centering
\caption{Effect size and detectability over seven seeds. Cohen's $d$ is
computed on the paired across-seed differences. \enquote{Seeds required} is
for 80\% power at the observed effect size and between-seed variance.}
\label{tab:power}
\begin{tabular}{lcrr}
\toprule
Effect & $\Delta$ QWK & Cohen's $d$ & Seeds required \\
\midrule
Ordinal and cost-sensitive head (RQ5) & $+0.0082$ & $\mathbf{1.97}$ & 3 \\
Full system vs baseline & $+0.0177$ & $\mathbf{1.91}$ & 3 \\
Typed heterogeneous edges (RQ1) & $+0.0093$ & $\mathbf{1.06}$ & 7 \\
\midrule
Disease-conditioned routing (RQ3) & $+0.0053$ & $0.39$ & 52 \\
Anatomical topology vs shuffle (RQ1) & $+0.0020$ & $0.20$ & 205 \\
Modality dropout & $-0.0026$ & $-0.13$ & 442 \\
Cross-sequence self-supervision (RQ2) & $+0.0020$ & $0.12$ & 545 \\
Gated residual & $+0.0014$ & $0.07$ & 1{,}615 \\
Relational message passing & $+0.0002$ & $0.01$ & 65{,}489 \\
\bottomrule
\end{tabular}
\end{table}

The table separates into two groups with nothing in between. Three effects are
large by any convention, $d$ between $1.06$ and $1.97$, and all three are the
comparisons that survive correction. Everything else is $d < 0.4$. That gap
matters: a study genuinely starved of power would show a spread of intermediate
effects rather than a clean division, and would not resolve any comparison at
three seeds.

The ladder is therefore not an instrument that reports nothing everywhere. Two
effects are detectable with a single seed and a third with seven, which is what
the campaign used. Where it reports nothing, that is a statement about the
effect rather than about the instrument.

The largest standardised effect in the study is the ordinal and cost-sensitive
head at $d = 1.97$, exceeding the full-system comparison it forms part of. Its
between-seed standard deviation is $0.0042$, the smallest of any comparison, so
its modest raw difference is also the most reliably reproduced.

These figures supersede a three-seed version of the same table, and the
differences are instructive. Three-seed variance estimates are themselves noisy:
routing moved from an estimated $3{,}467$ seeds to $52$, and cross-sequence
self-supervision from $216$ to $545$. The ordering is stable and the two groups
are unchanged, but individual requirements should be read as order-of-magnitude
guidance rather than as precise targets.

\subsection{What the Nulls Establish}
\label{sec:results-equivalence}

Reporting a comparison as \enquote{not significant} states only that an effect
was not detected. It is the weakest available form of a null, and it invites the
reader to supply the explanation that the study was too small. The paired
interval already bounds the effect, and stating that bound converts an absence
into a measurement.

Table~\ref{tab:bounds} gives, for each comparison that did not separate, the far
edge of its 95\% interval: the largest effect still compatible with the data.

These are \emph{confidence} bounds and are deliberately not called equivalence
bounds. A formal equivalence test requires a smallest effect of interest to be
specified in advance and tested against, for example by two one-sided tests.
No such margin was pre-specified here, and choosing one now, with the results
already seen, would be post hoc. The bounds below therefore say what the data
are compatible with; they do not assert equivalence to zero.

\begin{table}[htbp]
\centering
\caption{Upper compatibility bounds for the comparisons that did not
separate. Each figure is the far edge of the 95\% confidence interval: the
largest effect still compatible with the data, in QWK and as a percentage
of the E0 baseline of $0.7270$. These are confidence bounds, not the
outcome of an equivalence procedure.}
\label{tab:bounds}
\begin{tabular}{lrr}
\toprule
Comparison & Effect is at most & As \% of baseline \\
\midrule
Gated residual & $0.0064$ & $0.89\%$ \\
Relational message passing & $0.0074$ & $1.02\%$ \\
Anatomical topology vs shuffle (RQ1) & $0.0076$ & $1.05\%$ \\
Cross-sequence self-supervision (RQ2) & $0.0089$ & $1.22\%$ \\
Modality dropout & $0.0091$ & $1.25\%$ \\
Disease-conditioned routing (RQ3) & $0.0112$ & $1.55\%$ \\
\bottomrule
\end{tabular}
\end{table}

The claims available are therefore considerably stronger than the absence of a
significant result. Anatomically correct graph topology contributes at most
$1.05\%$ of baseline performance over a degree-preserving random shuffle.
Anatomically aligned cross-sequence self-supervision contributes at most
$1.22\%$ over the same architecture without it. Disease-conditioned routing
contributes at most $1.55\%$ over fixed fusion.

Set against the reference standard, these bounds are small in a way that matters.
Section~\ref{sec:lr-reliability} reports inter-reader $\kappa$ between $0.49$ and
$0.73$ depending on compartment. An effect bounded below $1.6\%$ of baseline is
not merely undetectable in this study; it is below the level at which the
reference standard itself could adjudicate the question.

\subsection{The Seven-Seed Extension}
\label{sec:results-sevenseed}

The campaign was extended from three seeds to seven, and the extension is
reported rather than absorbed, because it changed two verdicts.

Typed heterogeneous edges moved from $+0.0123$ on 3/3 seeds with
$p_{\mathrm{FDR}} = 0.054$ --- narrowly missing correction --- to $+0.0093$ on
7/7 seeds with $p_{\mathrm{FDR}} = 0.009$. The point estimate fell while the
between-seed standard deviation fell further, and the comparison crossed the
threshold. The ordinal head moved the same way, from $0.054$ to $0.009$.

Anatomical topology moved in the opposite direction: from $+0.0051$ on 2/3 seeds
to $+0.0020$ on 4/7. So did cross-sequence self-supervision, from $+0.0051$ on
2/3 to $+0.0020$ on 4/7.

That divergence is the reason the extension is worth reporting. The same four
additional runs that promoted two comparisons to significance made two others
weaker. Insufficient power is therefore not available as a blanket explanation
for the nulls in this chapter: the instrument demonstrably resolved effects at
this scale, and the effects it did not resolve moved away from significance as
evidence accumulated rather than toward it.

\subsection{Clinical Error Structure}
\label{sec:results-errors}

Aggregate agreement conceals the errors that matter clinically. A
Severe$\rightarrow$Normal/Mild confusion may withhold treatment from a patient
who needs it, and the cost matrix of Section~\ref{sec:method-cost} was chosen to
suppress that direction specifically.

\begin{table}[htbp]
\centering
\caption{Clinical error structure over seven seeds. \enquote{Distance $\geq 2$}
is the rate of predictions two or more grades from the reference.}
\label{tab:errors}
\begin{tabular}{lrrr}
\toprule
Configuration & Severe recall & Severe $\rightarrow$ Normal & Distance $\geq 2$ \\
\midrule
E0 single sequence & 61.1\% & 5.7\% & 0.526\% \\
E1 multi-sequence & 61.1\% & 4.9\% & 0.500\% \\
E2 + disease routing & 60.7\% & 4.7\% & 0.414\% \\
E3 + modality dropout & 61.4\% & 4.6\% & 0.420\% \\
E4 + ACSSL & 58.8\% & 4.0\% & 0.453\% \\
E5 + homogeneous graph & 59.0\% & 4.1\% & 0.389\% \\
E6 + typed graph & 62.9\% & 4.1\% & 0.375\% \\
\quad E6 shuffled edges & 59.3\% & 3.7\% & 0.330\% \\
\quad E6 ungated residual & \textbf{65.1\%} & 4.3\% & 0.500\% \\
E7 + ordinal, cost-sensitive & 63.1\% & \textbf{3.8\%} & \textbf{0.344\%} \\
\bottomrule
\end{tabular}
\end{table}

Severe$\rightarrow$Normal errors fall from $5.7\%$ at E0 to $3.8\%$ at E7, a
relative reduction of a third, while recall on Severe targets rises from
$61.1\%$ to $63.1\%$. Distant errors fall by a comparable margin, from $0.526\%$
to $0.344\%$. The system does not purchase aggregate agreement by trading away
the errors clinicians care most about; both move in the desired direction
together.

One complication should be stated rather than smoothed over. The ungated variant
achieves the highest Severe recall of any configuration at $65.1\%$, and
simultaneously carries the joint-highest rate of distant errors at $0.500\%$
against E7's $0.344\%$. Greater sensitivity bought with more badly misplaced
predictions is precisely the trade the cost matrix exists to refuse, and E7
makes it in the intended direction.

\subsection{Calibration}
\label{sec:results-calibration}

The expected calibration error recorded for each run is the \emph{uncalibrated}
test value; temperature scaling is fitted on the validation partition and its
test metrics are stored separately. The distinction is material, because the
uncalibrated figures alone would suggest the ladder becomes progressively worse
calibrated, and the correction the protocol actually applies reverses that.

\begin{table}[htbp]
\centering
\caption{Expected calibration error before and after temperature scaling, and
the fitted temperature, over seven seeds.}
\label{tab:calibration}
\begin{tabular}{lrrr}
\toprule
Configuration & ECE uncalibrated & ECE calibrated & Temperature \\
\midrule
E0 single sequence & $0.0293$ & $\mathbf{0.0130}$ & $1.161$ \\
E1 multi-sequence & $0.0214$ & $0.0129$ & $1.131$ \\
E2 + disease routing & $0.0428$ & $0.0168$ & $1.601$ \\
E3 + modality dropout & $0.0326$ & $0.0145$ & $1.200$ \\
E4 + ACSSL & $0.0266$ & $0.0110$ & $1.187$ \\
E5 + homogeneous graph & $0.1024$ & $0.0397$ & $2.019$ \\
E6 + typed graph & $0.0426$ & $0.0209$ & $1.422$ \\
E7 + ordinal, cost-sensitive & $0.0508$ & $0.0245$ & $1.306$ \\
\bottomrule
\end{tabular}
\end{table}

Temperature scaling roughly halves ECE at every rung, and every fitted
temperature exceeds $1$, so every configuration is overconfident before
correction.

Two observations run against the system and are reported for that reason. E5,
the homogeneous graph, is by a wide margin the worst calibrated configuration at
$0.1024$ uncalibrated with a fitted temperature of $2.019$; it is also the rung
that fails to improve accuracy, and its overconfidence is a second symptom of
the same failure. And the accumulating ladder does not improve calibration: E7
at $0.0245$ is worse than E0 at $0.0130$, so the gains in agreement come with
less reliable probabilities. RQ5 asked whether calibrated uncertainty could
support selective prediction. On this evidence the agreement half of that
question is answered affirmatively and the calibration half is not.

% ===========================================================================
\section{RQ1 --- Relational Disease Modelling}
\label{sec:results-rq1}

\begin{quote}
\emph{Does representing lumbar disease targets as a typed heterogeneous graph
improve grading over independent target heads, an ordered five-level transformer
and a homogeneous graph, and are any gains concentrated where single-target
evidence is weakest?}
\end{quote}

RQ1 divides into two claims, and the evidence separates them.

\subsection{Typed Structure Helps}

A typed heterogeneous graph outperforms a homogeneous one by $+0.0093$ QWK
($[+0.0027, +0.0156]$), and does so on \textbf{every one of seven seeds}. The
comparison survives false discovery rate correction at
$p_{\mathrm{FDR}} = 0.009$, with $d = 1.06$.

This is one of the two verdicts the seven-seed extension changed. At three seeds
the same comparison stood at $+0.0123$ with $p_{\mathrm{FDR}} = 0.054$, narrowly
outside correction, and the power analysis then indicated roughly ten seeds would
be needed. Seven sufficed, because the additional runs reduced the between-seed
standard deviation faster than they reduced the point estimate.

Relational message passing alone does not help: E5 against E0 is $+0.0002$
($[-0.0073, +0.0074]$, $p_{\mathrm{FDR}} = 0.967$), which is as close to exactly
nothing as this study measures. The benefit appears only once relations are
typed. Message passing over an untyped graph is, on this evidence, worth
precisely its parameter count and no more.

\subsection{Anatomical Topology Does Not}

H3 predicted that \enquote{performance under a shuffled-edge control should
deteriorate if the topology itself matters}. It does not deteriorate
meaningfully. E6 exceeds its shuffled control by $+0.0020$ QWK
($[-0.0038, +0.0076]$, 4/7 seeds, $p_{\mathrm{FDR}} = 0.688$), and the paired
interval bounds any advantage at $1.05\%$ of baseline performance. The control
preserves node count, edge count and degree, so the two models are matched in
capacity and differ only in which targets are connected to which.

The conclusion is therefore sharper than the hypothesis anticipated.
\emph{Typing the relations carries information; the anatomical identity of the
relations does not.} What the graph supplies is a structured parameterisation of
interaction between targets, not the specific adjacency of levels, compartments
and bilateral pairs that motivated it.

\subsection{The Gating Ablation}

Chapter~\ref{ch:methodology} designates the ungated residual a mandatory
ablation. The gated and ungated formulations are indistinguishable
($+0.0014$, 3/7 seeds, $p_{\mathrm{FDR}} = 0.713$), with any difference bounded
at $0.89\%$ of baseline --- the tightest bound in the study. The gate can be
removed without cost.

\paragraph{Answer to RQ1.} Partly. Typed relational structure improves grading
over both independent heads and a homogeneous graph. The anatomical content of
that topology is not what produces the improvement, and the gating mechanism
contributes nothing.

% ===========================================================================
\section{RQ2 --- Anatomically Aligned Self-Supervision}
\label{sec:results-rq2}

\begin{quote}
\emph{Does cross-sequence pretraining based on DICOM-defined anatomical
correspondence improve grading, label efficiency and transfer?}
\end{quote}

Pretraining itself succeeded as an optimisation: validation InfoNCE reached
$1.5162$ against a chance value of $3.4657$, so the encoders demonstrably learned
to match corresponding anatomy across sequences. The question is whether that
representation is worth anything downstream. Three independent probes say it is
not.

\emph{Accuracy.} E4 against E3 is $+0.0020$ QWK ($[-0.0043, +0.0089]$, 4/7
seeds, $p_{\mathrm{FDR}} = 0.688$, $d = 0.12$), with any benefit bounded at
$1.22\%$ of baseline. The seven-seed extension moved this comparison
\emph{away} from significance, not toward it: at three seeds it stood at
$+0.0051$ on 2/3 seeds.

\emph{Cross-sequence robustness.} This is the probe on which anatomical
self-supervision should be strongest, because a representation that has learned
one anatomy across three acquisitions ought to lose less when one acquisition is
withheld. Under the controlled input ablation of
Section~\ref{sec:results-rq3-ablation}, adding ACSSL changes reliance on the
annotated sequence by $-0.0069 \pm 0.0304$ on 2/3 seeds (this probe was run on
seeds 42--44 only) --- indistinguishable
from zero. By contrast, ordinary modality dropout reduces that reliance by
$-0.0588$ on 3/3 seeds.

\emph{Attention.} Under the attribution analysis of
Section~\ref{sec:results-attribution}, E4 is the only configuration whose
attention concentration falls \emph{below} the single-sequence baseline
($-0.027$ against E0, 1/3 seeds).

\paragraph{Answer to RQ2.} No. H1 is rejected on all three axes tested. The
pretext task is learnable and is learned; the representation it produces confers
no measurable downstream benefit at this data scale. Label efficiency and
external transfer were not separately evaluated
(Section~\ref{sec:results-rq4}), so the rejection is scoped to internal grading
and robustness.

% ===========================================================================
\section{RQ3 --- Disease-Conditioned Modality Routing}
\label{sec:results-rq3}

\begin{quote}
\emph{Can one model learn target- and patient-dependent sequence weights while
remaining robust to missing and quality-degraded MRI sequences?}
\end{quote}

\subsection{The Learned Allocation}

The router learns the clinically expected allocation with complete consistency.
Across all fifteen runs in which a gate is present, foraminal targets allocate
most weight to sagittal T1, central canal targets to sagittal T2, and
subarticular targets to axial T2. The pattern reproduced 15/15 without being
encoded as a prior.

Chapter~\ref{ch:methodology} is explicit that this is insufficient:
\enquote{the gate weights are model allocations, not causal estimates of sequence
importance}.

\subsection{Controlled Input Ablation}
\label{sec:results-rq3-ablation}

The prescribed intervention withholds each sequence at test time and measures the
per-condition loss. Withholding the sequence a condition is graded from is
severely damaging --- between $0.06$ and $0.35$ QWK, an order of magnitude larger
than any difference in Table~\ref{tab:comparisons} --- while withholding the
others is near-neutral or slightly beneficial. Agreement between allocation and
intervention is 5/5 conditions on every seed.

Taken alone, this reads as strong causal support for RQ3. It is not, and the
control is what reveals why.

\emph{E1 has no router.} It fuses the three sequences with fixed weights. It
produces the same 5/5 pattern with drops as large or larger. The dependence is a
property of the dataset: each condition is \emph{annotated} on one particular
sequence, so withholding that sequence removes the only crop actually centred on
the pathology. Any model suffers, routed or not.

Isolating each step's effect on mean reliance on the annotated sequence:

The input ablation was run on seeds 42--44, before the campaign was extended,
and its figures are three-seed.

\begin{center}
\begin{tabular}{lcc}
\toprule
Step & $\Delta$ reliance & Seeds lower \\
\midrule
Add the router             & $-0.0010 \pm 0.0382$ & 1/3 \\
Add modality dropout       & $\mathbf{-0.0588 \pm 0.0561}$ & \textbf{3/3} \\
Add ACSSL                  & $-0.0069 \pm 0.0304$ & 2/3 \\
\bottomrule
\end{tabular}
\end{center}

The router changes sequence reliance by essentially zero. Its gate weights track
the annotation structure of the dataset without creating any causal dependence
beyond what fixed fusion already possesses.

\paragraph{Answer to RQ3.} The first half is yes and the second half is no. One
model does learn target-dependent sequence weights, and those weights are
clinically sensible. But they do not improve grading ($+0.0053$ QWK, 5/7 seeds,
$p_{\mathrm{FDR}} = 0.131$, bounded at $1.55\%$ of baseline)
and they do not improve robustness to missing sequences. H2 predicted that
routing would \enquote{degrade less under missing-sequence experiments than fixed
concatenation}; it does not. The graceful degradation that the system does
exhibit is attributable to modality dropout, a standard regulariser that
Chapter~\ref{ch:introduction} explicitly does not claim as original.

% ===========================================================================
\section{RQ4 --- Cross-Institutional Transfer}
\label{sec:results-rq4}

\textbf{Not yet answered.} The zero-shot transfer experiment on the Rizgary
Hospital cohort has not been executed. Reporting its status honestly is required
by the protocol changelog commitment of
Section~\ref{sec:intro-hypotheses}.

Preparatory work is complete and is documented in
\texttt{thesis/chapter4/}: the cohort has been surveyed, radiology reports parsed
into structured labels, imaging and report records reconciled, and a grading
worklist produced. Of the cases with both a report and usable imaging, sixteen
carry conflicts between the report-derived label and the spreadsheet reference
that require adjudication by a reader before any transfer number can be
computed.

Two blockers stand between the present state and an RQ4 result. First, the
sixteen conflicting cases await regrading. Second, the cohort DICOMs carry
unredacted patient identifiers, and the existing de-identification routine keys
on a \texttt{PatientID} field that is not unique in this cohort --- $45$ distinct
values across $346$ cases --- so applying it would merge roughly eight patients
into each pseudonym. Neither blocker is technical in the sense of requiring new
method; both require decisions and reader time.

H4 and H5 are therefore untested. They are retained here rather than removed, in
keeping with the requirement that components not supported by executed
experiments be reported as such rather than deleted from the research history.

% ===========================================================================
\section{RQ5 --- Clinically Asymmetric Error}
\label{sec:results-rq5}

\begin{quote}
\emph{Do ordinal and cost-sensitive objectives reduce clinically consequential
distant errors, particularly Severe$\rightarrow$Normal/Mild?}
\end{quote}

Yes, and this is the second most reliable effect in the study. Adding the ordinal
and cost-sensitive head improves QWK by $+0.0082$ ($[+0.0025, +0.0143]$,
$p_{\mathrm{FDR}} = 0.009$, $d = 1.97$ --- the largest standardised effect in the
study) on 7/7
seeds, with the smallest between-seed standard deviation of any comparison
($0.0017$) --- so small that a single seed suffices to detect it
(Table~\ref{tab:power}).

The clinically material quantity moves in the intended direction. Recall on
Severe targets rises from $62.7\%$ at E0 to $65.0\%$ at E7, and from $59.6\%$ at
E6, the same architecture under categorical cross-entropy. The objective is
doing what it was selected to do rather than trading distant errors for
aggregate accuracy.

H6 anticipated that these objectives \enquote{may reduce distant errors even if
they do not improve aggregate macro-F1}. In the event both improved.

\paragraph{Answer to RQ5.} Yes. This is noted with some irony: the objective
function, explicitly disclaimed in
Section~\ref{sec:intro-contrib-supporting} as a supporting component and not an
original method, is a larger and far more reliable contributor to final
performance than two of the three proposed novel mechanisms.

% ===========================================================================
\section{Why the Structural Priors Had Little to Add}
\label{sec:results-attribution}

The results above establish that most of the proposed mechanisms do not improve
grading. A null result with a mechanism is worth more than a null result with a
p-value, and this section supplies the mechanism.

\subsection{Method}

Each crop is centred by construction on its target's annotated coordinate, so the
pathology lies at the middle of the frame and attention on it can be measured
rather than described: the fraction of Grad-CAM mass falling inside a central
disc of $15$~mm radius at the $60$~mm field of view.

That disc covers $19.7\%$ of the frame, which is what a uniform map would score.
Because convolutional networks are centre-biased irrespective of what they have
learned, the identical architecture with untrained weights is evaluated as a
floor; it scores $0.204$.

Attribution depth was fixed on resolution grounds before results were compared. A
residual network downsamples by $32$, so at these $128$~pixel crops the final
block is $4 \times 4$ --- one cell spanning $15$~mm, too coarse to resolve the
disc. The penultimate block at $8 \times 8$ matches the $7 \times 7$ resolution
at which Grad-CAM is conventionally applied to $224$~pixel inputs, and is used
throughout. All three depths are reported in the appendix so that the dependence
is visible rather than taken on trust.

\subsection{Result}

\begin{center}
\begin{tabular}{lccc}
\toprule
Configuration & CAM mass in disc & vs.\ untrained & vs.\ E0 \\
\midrule
E0 single sequence     & $0.490 \pm 0.058$ & $+0.286$ & --- \\
E1 multi-sequence      & $0.578 \pm 0.073$ & $+0.374$ & $+0.089$ (3/3) \\
E2 + routing           & $0.580 \pm 0.063$ & $+0.376$ & $+0.090$ (2/3) \\
E3 + modality dropout  & $0.515 \pm 0.071$ & $+0.311$ & $+0.026$ (2/3) \\
E4 + ACSSL             & $0.462 \pm 0.027$ & $+0.258$ & $-0.027$ (1/3) \\
\midrule
Untrained (same architecture) & $0.204$ & --- & --- \\
Uniform map (disc area)       & $0.197$ & --- & --- \\
\bottomrule
\end{tabular}
\end{center}

Every trained configuration places between $2.3$ and $2.9$ times the chance share
of its attention on the annotated centre. Critically, \textbf{E0 already reaches
$0.490$} --- a plain single-sequence convolutional network with no routing, no
self-supervision and no graph.

This is the explanation the null results require. The structural priors of RQ1--RQ3
were motivated by the premise that anatomical context is needed to find and
interpret the lesion. The encoder locates it without them. There is little left
for a prior to contribute, and the ladder measures exactly that little.

\subsection{A Dissociation Worth Noting}

Multi-sequence input is the only step that improves attention on all three seeds
($+0.089$ against E0), and it is simultaneously the step with the most negative
accuracy effect ($-0.0055$ QWK). Additional sequences sharpen \emph{where} the
model looks without improving \emph{what it concludes}. Attention concentration
and grading accuracy are not the same quantity, and this study can separate them.

% ===========================================================================
\section{Threats to Validity}
\label{sec:results-threats}

\subsection{Errors Detected in the Analysis}

Three analyses produced incorrect conclusions before controls were added. Each is
recorded because each would have been reported as a positive finding, and because
all three erred in the direction of the hypothesis under test.

\emph{Per-seed intervals reversed sign.} Described in
Section~\ref{sec:results-inference}. Two comparisons yielded opposite,
individually significant results on different seeds.

\emph{A mechanism check without a control read as proof.} The 5/5 agreement
between routing weights and intervention (Section~\ref{sec:results-rq3-ablation})
appeared to establish RQ3 causally, with effects an order of magnitude larger
than any ladder difference, until a router-free model produced the same pattern.

\emph{An attribution probe measured the wrong encoder.} The first
implementation hooked the first encoder unconditionally, which for multi-sequence
configurations is the sagittal T1 encoder --- including on subarticular targets
graded from axial T2. It reported that no configuration improved on E0
($-0.036$, 1/3 seeds). Selecting the encoder that read each target's annotated
sequence reverses the sign to $+0.089$ on 3/3 seeds.

\subsection{Limitations}

\emph{Seed count.} Three seeds is thin for effects of this magnitude; a
seven-seed campaign is running. Table~\ref{tab:power} indicates that seven seeds
will resolve typed edges but will not resolve RQ2 or RQ3, whose requirements are
$216$ and $3{,}467$ seeds respectively.

\emph{Single benchmark.} All results are internal to LumbarDISC. RQ4 is
unanswered, so external validity is untested and no claim of generalisation is
made.

\emph{Point annotations.} Centre concentration presumes the annotated coordinate
marks the lesion. RSNA coordinates come from single readers; a systematic offset
would depress the measure uniformly across configurations.

\emph{Attribution method.} Grad-CAM is one attribution technique with known
failure modes. The untrained floor makes the comparison meaningful, but absolute
values are method-dependent.

\emph{The input ablation covers E0--E4 only.} Those rungs reach the fusion stage
directly, whereas E5--E7 route through the graph. The routing claim of RQ3 is in
any case cleanest to test at E2, with the graph out of the path.

\emph{Attribution on the graph rungs measures a different quantity for E7.}
Attribution was extended to E5--E7, which required one backward pass per node
rather than one per batch: message passing couples the nodes, so a summed
backward pass reports what an image contributed to the \emph{whole graph} rather
than to its own target, understating concentration by $0.173$ on the nodes
examined. With that corrected, E5 and E6 reach $0.341$ and $0.428$ against
untrained floors near $0.18$.

E7 pools to $0.171$, below chance, and that figure is a prevalence artefact
rather than a property of the model. Its ordinal head explains an accumulated
severity score, so asking what makes a target look severe highlights the
\emph{absence} of evidence on a target with no pathology, and $77.33\%$ of
targets are Normal/Mild. Split by grade, E7 gives $0.098$, $0.269$ and $0.415$
for Normal/Mild, Moderate and Severe, while E6 is flat at $0.429$, $0.419$ and
$0.448$. On targets that carry pathology the full system localises as well as
the rung below it. E7's pooled value is therefore not comparable with the other
rungs and is not presented as though it were.

% ===========================================================================
\section{Summary of Answers}
\label{sec:results-summary}

\begin{center}
\begin{tabular}{llp{6.6cm}}
\toprule
RQ & Verdict & Basis \\
\midrule
RQ1 & Partly & Typed edges $+0.0093$, \textbf{7/7 seeds},
$p_{\mathrm{FDR}} = 0.009$. Anatomical topology does not beat a
degree-preserving shuffle ($+0.0020$, 4/7), bounded at $1.05\%$ of baseline. \\
RQ2 & No & $+0.0020$, 4/7 seeds, bounded at $1.22\%$; no robustness benefit;
attention below baseline. Rejected on three axes, and weakened by the
seven-seed extension. \\
RQ3 & Mechanism yes, benefit no & Allocation replicates 15/15 but a router-free
model matches it under intervention. $+0.0053$ QWK, 5/7 seeds, bounded at
$1.55\%$. \\
RQ4 & Not executed & Cohort preparation complete; blocked on reader
adjudication, de-identification, and the absence of localisation coordinates. \\
RQ5 & Yes & $+0.0082$, \textbf{7/7 seeds}, $p_{\mathrm{FDR}} = 0.009$, $d = 1.97$
--- the largest standardised effect in the study. Severe$\rightarrow$Normal
errors fall $5.7\% \rightarrow 3.8\%$. \\
\midrule
System & Yes & $+0.0177$ QWK $[+0.0097, +0.0260]$, \textbf{7/7 seeds},
$p_{\mathrm{FDR}} < 0.001$. \\
\bottomrule
\end{tabular}
\end{center}

The thesis set out to determine whether anatomical structure, imposed explicitly
on a multi-sequence grading model, improves degenerative lumbar classification.
The answer this chapter supports is that it largely does not, that the
convolutional encoder already performs the localisation such structure was meant
to supply, and that the components which do improve performance are relational
typing and a clinically motivated objective rather than anatomical correspondence
or disease-conditioned routing.

Two qualifications belong with that summary. The improvement is bounded by a
reference standard on which expert readers agree at $\kappa$ between $0.49$ and
$0.73$, so its ceiling is set by the labels rather than by the architecture. And
the accumulating ladder improves agreement without improving the reliability of
its probabilities: calibrated ECE is worse at E7 than at E0.

Chapter~\ref{ch:discussion} considers what follows for the design of anatomy-aware
medical imaging models, and why mechanisms that are individually well motivated
and demonstrably learnable can nonetheless fail to pay.

\end{document}
